% !TEX root = msfp26-lists.tex
\documentclass[acmsmall,natbib=false,nonacm,review,screen]{acmart}
\settopmatter{printfolios=false,printccs=false,printacmref=false}
\setcopyright{none}

\AtBeginDocument{%
  \providecommand\BibTeX{{Bib\TeX}}}

\RequirePackage[
  datamodel=acmdatamodel,
  style=acmauthoryear,
]{biblatex}

\addbibresource{references.bib}

\usepackage[appendix=append,bibengine=biblatex]{apxproof}

\newtheoremrep{theorem}{Theorem}

\usepackage{subcaption}
\usepackage{enumitem}

\usepackage{xcolor-solarized}
\colorlet{magenta}{solarized-magenta}
\colorlet{cyan}{solarized-cyan}
\colorlet{blue}{solarized-blue}
\colorlet{green}{solarized-green}
\colorlet{nilbg}{solarized-yellow!25}
\colorlet{consbg}{solarized-violet!25}
\colorlet{Appbg}{solarized-blue!25}
\colorlet{nbg}{solarized-cyan!25}
\colorlet{cbg}{solarized-red!25}
\colorlet{itbg}{solarized-blue!25}

\usepackage{rv-macros}

\begin{document}

\title{Symmetric List Objects}

\author{Fiona Blackett}
\email{fiona.blackett@strath.ac.uk}
\affiliation{%
  \institution{University of Strathclyde}
  \country{UK}}

\author{Vikraman Choudhury}
\email{me@vikraman.org}
\affiliation{%
  \institution{University of Strathclyde}
  \country{UK}}

\author{Rin Liu}
\email{rin.liu@strath.ac.uk}
\affiliation{%
  \institution{University of Strathclyde}
  \country{UK}}

\maketitle

\vspace*{-2ex}

\section{Introduction}
\label{sec:introduction}

Lists are a fundamental data structure, important in computer science and mathematics. They are often defined in functional programming languages as cons-lists, generated inductively from a pair of constructors $\nil:L(A)$ and $\cons:A\times L(A)\to L(A)$, and whose eliminators allow for definitions by recursion.

In category theory, these lists are understood as initial algebras $\mu X.1+A\times X$~\cite{,DBLP:phd/ethos/Hagino87,adamekInitialAlgebrasTerminal2025}. This structure can be axiomatised more generally in the setting of a monoidal category, replacing the cartesian product with a tensor, and removing the need for a coproduct by axiomatising the $\nil$ and $\cons$ maps directly with an iteration mechanism. These list objects generalise the notion of Natural Numbers Objects~\cite{lawvereFUNCTORIALSEMANTICSALGEBRAIC1963} and were introduced in Joyal's arithmetic universes. An internal language for them was studied by Maietti~\cite{maiettiJoyalsArithmeticUniverse2010}, and other variants were explored in~\cite{cockettListarithmeticDistributiveCategories1990} and~\cite{jayFiniteObjectsLocos1997}.

In all of these settings, additional axioms on the underlying category are required to reconstruct traditional list operations from these list objects. However, in the minimal setting of a monoidal category only, this notion of a list object can be strengthened to that of a \textit{parametrised} list object which supports parametrised recursion. It is shown in~\cite{fioreListObjectsAlgebraic2017} that this yields a construction of free monoids (see also~\cite{kellyUnifiedTreatmentTransfinite1980}) in monoidal categories. We give a string diagrammatic proof of this result, showing that if a monoidal category $\cat{C}$ has all parametrised list objects, then the forgetful functor $\mathrm{Mon}(\cat{C})\to\cat{C}$ is monadic.

Finite multisets are another important data structure that are the commutative counterpart to lists; and in fact can be axiomatised with only a simple modification to these list objects. In this talk, we will study the notion of a \textit{symmetric} list object in a symmetric monoidal category. We will identify lists that differ by a transposition of the two frontmost elements using the symmetry of the underlying category. Using this, we show the analogous commutative results that symmetric list objects are free commutative monoids, and that having all symmetric list objects yields a similar monadic adjunction between $\mathrm{CMon}(\cat{C})$ and $\cat{C}$.

In type theory, finite multisets have been considered as setoids over lists, where the equivalence relation is bag equivalence~\cite{danielssonBagEquivalenceProofRelevant2012}, or as higher inductive types of lists with an additional swap constructor to permute the two head elements~\cite{choudhuryFreeCommutativeMonoids2023}. There, the ambient category is cartesian closed, but our setting only requires a symmetric monoidal structure.

\section{List objects}
\label{sec:list-objects}

Let $\cat{C}(1, \otimes)$ be a strict monoidal category.

\begin{definition}
  A \textit{parametrised list object} (hereafter, just \textit{list object}) $L(A)$ on an object $A$ consists of a pair of maps $\nil:1\to L(A)$ and $\cons:A\otimes L(A)\to L(A)$
  % \[
  %   \Nullary{\nil}{L(A)} \quad \text{and} \quad \smash{\Binary{A}{L(A)}{\cons}{L(A)}}
  % \]
  such that for every object $L$ equipped with maps $n:P\to L$ and $c:A\otimes L\to L$,
  % \[
  %   \Unary{P}{n}{L} \quad \text{and} \quad \smash{\Binary{A}{L}{c}{L}}
  % \]
  there exists a unique \textit{mediating map} or \textit{iterator} $\iter(n,c):L(A)\otimes P\to L$
  % \[
  %   \smash{\Binary{L(A)}{P}{\iter(n,c)}{L}}
  % \]
  satisfying computation rules (a) and (b):

  \begin{figure}[H]
    \centering
    \begin{subfigure}[t]{0.39\linewidth}
      \centering
      \begin{tikzpicture}[baseline=(op.base)]
        \begin{scope}
          \node[operator,fill=itbg] (op) at (0,0) {$\iter(n,c)$};
          \node[anchor=east,operator,fill=nilbg] (in1) at ($(op.140)+(-1,0.4)$) {$\nil$};
          \node[anchor=east] (in2) at ($(op.220)+(-1,-0.4)$) {};
          \node[anchor=west] (out) at ($(op.east)+(1,0)$) {};
          \draw[blue] (in1.east) to[out=0,in=140] node[inputwire,pos=0.4,above,sloped=false] {$L(A)$} (op.140);
          \draw[green] ($(in2.west) - (0.25,0)$) to[out=0,in=220] node[inputwire,below,sloped=false] {$P$} (op.220);
          \draw[cyan] (op.east) -- (out.west) node[outputwire,above] {$L$};
        \end{scope}

        \node at (2,0) {$=$};

        \begin{scope}[xshift=3.65cm]
          \node[operator,fill=nbg] (op) at (0,0) {$n$};
          \node[anchor=east] (in) at ($(op.west)+(-1,0)$) {};
          \node[anchor=west] (out) at ($(op.east)+(1,0)$) {};
          \draw[green] (in.east) -- node[inputwire,above] {$P$} (op.west);
          \draw[cyan] (op.east) -- node[outputwire,above] {$L$} (out.west);
        \end{scope}
      \end{tikzpicture}%
      \caption{$\nil$-$\beta$}
    \end{subfigure}
    \hfill
    \begin{subfigure}[t]{0.59\linewidth}
      \centering
      \begin{tikzpicture}[baseline=(it.base)]
        \begin{scope}
          \node[operator,fill=itbg] (it) at (0,0) {$\iter(n,c)$};
          \node[anchor=east,operator,fill=consbg] (cons) at ($(it.140)+(-1,0.4)$) {$\cons$};
          \node[anchor=east] (in3) at ($(it.220)+(-1,-0.6)$) {};
          \node[anchor=west] (out) at ($(it.east)+(1,0)$) {};
          \node[anchor=east] (in1) at ($(cons.140)+(-1,0.2)$) {};
          \node[anchor=east] (in2) at ($(cons.220)+(-1,-0.2)$) {};

          \draw[blue] (cons.east) to[out=0,in=140] node[inputwire,pos=0.4,above,sloped=false] {$L(A)$} (it.140);
          \draw[green] let \p1 = (in1.east), \p2 = (in3) in
            (\x1,\y2) to[out=0,in=220] node[inputwire,pos=0.1,below,sloped=false] {$P$} (it.220);

          \draw[cyan] (it.east) -- (out.west) node[outputwire,above] {$L$};
          \draw[magenta] (in1.east) to[out=0,in=140] node[inputwire,above,sloped=false] {$A$} (cons.140);
          \draw[blue] (in2.east) to[out=0,in=220] node[inputwire,below,sloped=false] {$L(A)$} (cons.220);
        \end{scope}

        \node at (2,0) {$=$};

        \begin{scope}[xshift=5.7cm]
          \node[operator,fill=cbg] (it) at (0,0) {$c$};
          \node[anchor=east,operator,fill=itbg] (cons) at ($(it.220)+(-1,-0.4)$) {$\iter(n,c)$};
          \node[anchor=east] (in1) at ($(it.140)+(-1,0.8)$) {};
          \node[anchor=west] (out) at ($(it.east)+(1.3,0)$) {};
          \node[anchor=east] (in2) at ($(cons.140)+(-1,0.2)$) {};
          \node[anchor=east] (in3) at ($(cons.220)+(-1,-0.2)$) {};

          \draw[magenta] (cons.east) to[out=0,in=220] node[inputwire,pos=0.4,below,sloped=false] {$A$} (it.220);
          \draw[magenta] let \p1 = (in2.east), \p2 = (in1) in
            (\x1,\y2) to[out=0,in=140] node[inputwire,pos=0.1,above,sloped=false] {$A$} (it.140);

          \draw[cyan] (it.east) -- (out.west) node[outputwire,above] {$L$};
          \draw[blue] (in2.east) to[out=0,in=140] node[inputwire,above,sloped=false] {$L(A)$} (cons.140);
          \draw[green] (in3.east) to[out=0,in=220] node[inputwire,below,sloped=false] {$P$} (cons.220);
        \end{scope}
      \end{tikzpicture}%
      \caption{$\cons$-$\beta$}
    \end{subfigure}
  \end{figure}
\end{definition}

% \begin{definition}
%   A \textit{monoid object} in $\cat{C}$ is an object $M$ equipped with a pair of maps $\varepsilon:1\to M$ and $\mu:M\otimes M\to M$
%   % \[
%   %   \Nullary{\varepsilon}{M} \quad \text{and} \quad \Binary{M}{M}{\mu}{M}
%   % \]
%   satisfying:
%   \begin{figure}[H]
%     \centering
%     \begin{subfigure}[t]{0.49\linewidth}
%       \centering
%       \begin{tikzpicture}
%         \begin{scope}
%           \node[operator] (op) at (0,0) {$\mu$};
%           \node[anchor=east,operator] (in1) at ($(op.140)+(-1,0.4)$) {$\varepsilon$};
%           \node[anchor=east] (in2) at ($(op.220)+(-1,-0.4)$) {};
%           \node[anchor=west] (out) at ($(op.east)+(1,0)$) {};
%           \draw (in1.east) to[out=0,in=140] node[inputwire,pos=0.4,above,sloped=false] {$M$} (op.140);
%           \draw ($(in2.west) - (0.1,0)$) to[out=0,in=220] node[inputwire,below,sloped=false] {$M$} (op.220);
%           \draw (op.east) -- (out.west) node[outputwire,above] {$M$};
%         \end{scope}

%         \node at (2,0) {$=$};

%         \begin{scope}[xshift=3cm]
%           \draw (-0.5,0) -- node[midway,above] {$M$} (0.5,0);
%         \end{scope}

%         \node at (4,0) {$=$};

%         \begin{scope}[xshift=6cm]
%           \node[operator] (op) at (0,0) {$\mu$};
%           \node[anchor=east] (in1) at ($(op.140)+(-1,0.4)$) {};
%           \node[anchor=east,operator] (in2) at ($(op.220)+(-1,-0.4)$) {$\varepsilon$};
%           \node[anchor=west] (out) at ($(op.east)+(1,0)$) {};
%           \draw ($(in1.west) - (0.1,0)$) to[out=0,in=140] node[inputwire,above,sloped=false] {$M$} (op.140);
%           \draw (in2.east) to[out=0,in=220] node[inputwire,below,pos=0.4,sloped=false] {$M$} (op.220);
%           \draw (op.east) -- (out.west) node[outputwire,above] {$M$};
%         \end{scope}
%       \end{tikzpicture}%
%     \end{subfigure}
%     \hfill
%     \begin{subfigure}[t]{0.49\linewidth}
%       \centering
%       \begin{tikzpicture}
%         \begin{scope}
%           \node[operator] (op2) at (0,0) {$\mu$};
%           \node[anchor=east,operator] (op1) at ($(op2.140)+(-1,0.4)$) {$\mu$};
%           \node[anchor=west] (out) at ($(op2.east)+(1,0)$) {};
%           \node[anchor=east] (in1) at ($(op1.140)+(-1,0.2)$) {};
%           \node[anchor=east] (in2) at ($(op1.220)+(-1,-0.2)$) {};
%           \node[anchor=east] (in3) at ($(op2.220)+(-1,-0.6)$) {};

%           \draw (op1.east) to[out=0,in=140] node[inputwire,pos=0.4,above,sloped=false] {$M$} (op2.140);
%           \draw let \p1 = (in1.east), \p2 = (in3) in
%             (\x1,\y2) to[out=0,in=220] node[inputwire,pos=0.1,below,sloped=false] {$M$} (op2.220);

%           \draw (op2.east) -- (out.west) node[outputwire,above] {$M$};
%           \draw (in1.east) to[out=0,in=140] node[inputwire,above,sloped=false] {$M$} (op1.140);
%           \draw (in2.east) to[out=0,in=220] node[inputwire,below,sloped=false] {$M$} (op1.220);
%         \end{scope}

%         \node at (1.5,0) {$=$};

%         \begin{scope}[xshift=4.5cm]
%           \node[operator] (op1) at (0,0) {$\mu$};
%           \node[anchor=east,operator] (op2) at ($(op1.220)+(-1,-0.4)$) {$\mu$};
%           \node[anchor=west] (out) at ($(op1.east)+(1,0)$) {};
%           \node[anchor=east] (in1) at ($(op1.140)+(-1,0.6)$) {};
%           \node[anchor=east] (in2) at ($(op2.140)+(-1,0.2)$) {};
%           \node[anchor=east] (in3) at ($(op2.220)+(-1,-0.2)$) {};

%           \draw (op2.east) to[out=0,in=220] node[inputwire,pos=0.4,below,sloped=false] {$M$} (op1.220);
%           \draw let \p1 = (in2.east), \p2 = (in1) in
%             (\x1,\y2) to[out=0,in=140] node[inputwire,pos=0.1,above,sloped=false] {$M$} (op1.140);

%           \draw (op1.east) -- (out.west) node[outputwire,above] {$M$};
%           \draw (in2.east) to[out=0,in=140] node[inputwire,above,sloped=false] {$M$} (op2.140);
%           \draw (in3.east) to[out=0,in=220] node[inputwire,below,sloped=false] {$M$} (op2.220);
%         \end{scope}
%       \end{tikzpicture}%
%     \end{subfigure}
%   \end{figure}
% \end{definition}

\begin{theoremrep}
    The list object $L(A)$ is the free monoid on $A$.
\end{theoremrep}
\begin{proof}
    Define $\App:=\iter(\id,\cons):L(A)\otimes L(A)\to L(A)$. We claim that $(L(A),\nil,\App)$ is a free monoid structure on $A$.
    % We verify ($i$) left unit, ($ii$) right unit, ($iii$) associativity, and ($iv$) unique extensions.

    The left unit law directly holds by $\App$-$\nil$-$\beta$:
    \begin{equation*}
        \begin{tikzpicture}
            \begin{scope}
                \node[operator] (app) at (0,0) {$\App$};
                \node[anchor=east,operator] (nil) at ($(app.140)+(-1,0.4)$) {$\nil$};

                \node (in) at ($(app.220)+(-1.5,-0.4)$) {};
                \node (out) at ($(app.east)+(1,0)$) {};

                \draw[blue] (nil.east) to[out=0,in=140] (app.140);
                \draw[blue] (in) to[out=0,in=220]  (app.220);
                \draw[blue] (app.east) -- (out);
            \end{scope}

            \node at (2,0) {$=$};

            \begin{scope}[xshift=4cm]
                \draw[blue] (-1.25,0) -- (1.25,0);
            \end{scope}
        \end{tikzpicture}%
    \end{equation*}
    For the right unit law, consider the following maps $f,g:L(A)\to L(A)$ below:
    \begin{equation*}
        \begin{tikzpicture}
            \begin{scope}
                \node[operator] (app) at (0,0) {$\App$};
                \node[anchor=east,operator] (nil) at ($(app.220)+(-1,-0.4)$) {$\nil$};

                \node (in) at ($(app.140)+(-1.5,0.4)$) {};
                \node (out) at ($(app.east)+(1,0)$) {};

                \draw[blue] (in) to[out=0,in=140] node[inputwire,sloped=false] {$L(A)$} (app.140);
                \draw[blue] (nil.east) to[out=0,in=220] node[inputwire,below,pos=0.4,sloped=false] {$L(A)$} (app.220);
                \draw[blue] (app.east) -- (out) node[outputwire] {$L(A)$};
            \end{scope}

            \begin{scope}[xshift=4.5cm]
                \draw[blue] (-1.25,0) -- node[inputwire] {$L(A)$} (0,0);
                \draw[blue] (0,0) -- node[outputwire] {$L(A)$} (1.25,0);
            \end{scope}
        \end{tikzpicture}
    \end{equation*}
    We claim that these maps are both the iterator $\iter(\nil,\cons)$. By the uniqueness of the mediating map, it suffices to show that the two maps satisfy the required computation rules.
    
    The identity map trivially satisfies both rules, so we have $g=\iter(\nil,\cons)$. We have $f$-$\nil$-$\beta$ by applying $\App$-$\nil$-$\beta$:
    \begin{equation*}
        \begin{tikzpicture}
            \begin{scope}[xshift=-4.5cm]
                \node[operator,inner sep=1pt] (f) at (0,0) {$f$};
                \node[operator] (nil) at (-1.5,0) {$\nil$};

                \node (out) at ($(f.east)+(1,0)$) {};

                \draw[blue] (nil) -- node[inputwire] {} (f);
                \draw[blue] (f) -- node[inputwire] {} (out);
            \end{scope}
            
            \node at (-2.5,0) {$=$};

            \begin{scope}
                \node[operator] (op) at (0,0) {$\App$};
                \node[anchor=east,operator] (in1) at ($(op.140)+(-1,0.4)$) {$\nil$};
                \node[anchor=east,operator] (in2) at ($(op.220)+(-1,-0.4)$) {$\nil$};
                \node (out) at ($(op.east)+(1,0)$) {};
                \draw[blue] (in1.east) to[out=0,in=140] (op.140);
                \draw[blue] (in2.east) to[out=0,in=220] (op.220);
                \draw[blue] (op.east) -- (out);

                \begin{scope}[on background layer]
                    \node[
                        fit=(in1)(op),
                        diff,
                        rotate=-20,
                        inner sep=-3pt,
                    ] {};
                \end{scope}
            \end{scope}

            \node at (2,0) {$=$};

            \begin{scope}[xshift=4.5cm]
                \node[operator] (op) at (0,0) {$\id$};
                \node[anchor=east,operator] (in2) at ($(op.220)+(-1,-0.4)$) {$\nil$};
                \node (out) at ($(op.east)+(1,0)$) {};
                \draw[blue] (in2.east) to[out=0,in=220] (op.220);
                \draw[blue] (op.east) -- (out);
                \begin{scope}[on background layer]
                    \node[
                        fit=(op),
                        result,
                        inner sep=1.5pt,
                    ] {};
                \end{scope}
            \end{scope}

            \node at (6.5,0) {$=$};

            \begin{scope}[xshift=8.5cm]
                \node[operator] (op) at (-1,0) {$\nil$};
                \node (out) at ($(op.east)+(1,0)$) {};
                \draw[blue] (op.east) -- (out);
            \end{scope}
        \end{tikzpicture}%
    \end{equation*}
    and similarly, $f$-$\cons$-$\beta$ follows from $\App$-$\cons$-$\beta$:
    \begin{equation*}
        \begin{tikzpicture}
            \begin{scope}[xshift=-5.5cm]
                \node[operator,inner sep=1pt] (f) at (0,0) {$f$};
                \node[operator] (cons) at (-1.5,0) {$\cons$};

                \node (in1) at ($(cons.140)+(-1,0.2)$) {};
                \node (in2) at ($(cons.220)+(-1,-0.2)$) {};
                \node (out) at ($(f.east)+(0.8,0)$) {};

                \draw[magenta] (in1) to[out=0,in=140] (cons.140);
                \draw[blue] (in2) to[out=0,in=220]  (cons.220);
                \draw[blue] (cons) -- node[above] {$L(A)$} (f);
                \draw[blue] (f) -- (out);
            \end{scope}
            
            \node at (-3.8,0) {$=$};

            \begin{scope}
                \node[operator] (app) at (0,0) {$\App$};
                \node[anchor=east,operator] (cons) at ($(app.140)+(-1,0.4)$) {$\cons$};
                \node[anchor=east,operator] (nil) at ($(app.220)+(-2.2,-0.8)$) {$\nil$};
                
                \node (in1) at ($(cons.140)+(-1,0.2)$) {};
                \node (in2) at ($(cons.220)+(-1,-0.2)$) {};
                \node (out) at ($(app.east)+(0.8,0)$) {};

                \draw[magenta] (in1.east) to[out=0,in=140] (cons.140);
                \draw[blue] (in2.east) to[out=0,in=220] (cons.220);
                \draw[blue] (cons.east) to[out=0,in=140] (app.140);
                \draw[blue] (nil.east) to[out=0,in=220] (app.220);
                \draw[blue] (app.east) -- (out);
                
                \begin{scope}[on background layer]
                    \node[
                        fit=(cons)(app),
                        diff,
                        rotate=-20,
                        inner sep=-3pt,
                    ] {};
                \end{scope}
            \end{scope}

            \node at (1.8,0) {$=$};

            \begin{scope}[xshift=5.5cm]
                \node[operator] (cons) at (0,0) {$\cons$};
                \node[anchor=east,operator] (app) at ($(cons.220)+(-1,-0.4)$) {$\App$};
                \node[anchor=east,operator] (nil) at ($(app.220)+(-0.6,-0.2)$) {$\nil$};

                \node (in2) at ($(app.140)+(-1.2,0.2)$) {};
                \node (in1) at ($(in2 |- cons.140)+(0,0.6)$) {};
                \node (out) at ($(cons.east)+(0.8,0)$) {};

                \draw[magenta] (in1) to[out=0,in=140](cons.140);
                \draw[blue] (in2) to[out=0,in=140] (app.140);
                \draw[blue] (nil) to[out=0,in=220] (app.220);
                \draw[blue] (app) to[out=0,in=220] (cons.220);
                \draw[blue] (cons) -- (out);

                \begin{scope}[on background layer]
                    \node[
                        fit=(cons)(app),
                        result,
                        rotate=20,
                        inner sep=-3pt,
                    ] {};
                \end{scope}
            \end{scope}

            \node at (7.3,0) {$=$};

            \begin{scope}[xshift=10.2cm]
                \node[operator] (cons) at (0,0) {$\cons$};
                \node[operator] (f) at ($(cons.220)+(-1,-0.4)$) {$f$};

                \node (in2) at ($(f.west)+(-0.8,0)$) {};
                \node (in1) at ($(in2 |- cons.140) + (0,0.4)$) {};
                \node (out) at ($(cons.east)+(0.8,0)$) {};

                \draw[magenta] (in1) to[out=0,in=140] (cons.140);
                \draw[blue] (in2) -- (f);
                \draw[blue] (f) to[out=0,in=220] (cons.220);
                \draw[blue] (cons) -- (out);
            \end{scope}
        \end{tikzpicture}
    \end{equation*}
    and hence $f=\iter(\nil,\cons)=g$.
    
    For associativity, let $f,g:L(A)\otimes L(A)\otimes L(A)\to L(A)$ be the two maps below:
    \begin{equation*}
        \begin{tikzpicture}
            \begin{scope}
                \node[operator,fill=Appbg] (app1) at (0,0) {$\App$};
                \node[anchor=east,operator,fill=Appbg] (app2) at ($(app1.220)+(-1,-0.4)$) {$\App$};

                \node (in2) at ($(app2.140)+(-1,0.2)$) {};
                \node (in1) at ($(in2 |- app1.140)+(0,0.6)$) {};
                \node (in3) at ($(app2.220)+(-1,-0.2)$) {};
                \node (out) at ($(app1.east)+(1,0)$) {};

                \draw[blue] (in1) to[out=0,in=140] node[inputwire,pos=0.1,sloped=false] {$L(A)$} (app1.140);
                \draw[blue] (in2) to[out=0,in=140] node[inputwire,sloped=false] {$L(A)$} (app2.140);
                \draw[blue] (in3) to[out=0,in=220] node[inputwire,below,sloped=false] {$L(A)$} (app2.220);
                \draw[blue] (app1) -- (out) node[outputwire] {$L(A)$};
                \draw[blue] (app2) to[out=0,in=220] node[inputwire,pos=0.5,yshift=-2pt,below,sloped=false] {$L(A)$} (app1.220);
                
                \node (fill) at ($(app1.220)+(+1,-0.4)$) {};
            \end{scope}
            \begin{scope}[xshift=7.5cm]
                \node[operator,fill=Appbg] (app1) at (0,0) {$\App$};
                \node[anchor=east,operator,fill=Appbg] (app2) at ($(app1.140)+(-1,0.4)$) {$\App$};

                \node (in1) at ($(app2.140)+(-1,0.2)$) {};
                \node (in2) at ($(app2.220)+(-1,-0.2)$) {};
                \node (in3) at ($(in1 |- app1.220)+(0,-0.6)$) {};
                \node (out) at ($(app1.east)+(1,0)$) {};

                \draw[blue] (in1) to[out=0,in=140] node[inputwire,sloped=false] {$L(A)$} (app2.140);
                \draw[blue] (in2) to[out=0,in=220] node[inputwire,below,sloped=false] {$L(A)$} (app2.220);
                \draw[blue] (in3) to[out=0,in=220] node[inputwire,pos=0.1,below,sloped=false] {$L(A)$} (app1.220);
                \draw[blue] (app1) -- (out) node[outputwire] {$L(A)$};
                \draw[blue] (app2) to[out=0,in=140] node[inputwire,pos=0.5,yshift=2pt,sloped=false] {$L(A)$} (app1.140);

            \end{scope}
        \end{tikzpicture}
    \end{equation*}
    We claim that both maps are equal to $\iter(\App,\cons)$. We have $\nil$-$\beta$ for both since $\App$ computes on $\nil$ in the left argument. In detail, for $f$:
    \begin{equation*}
        \begin{tikzpicture}
            \begin{scope}
                \node[operator] (app2) at (0,0) {$\App$};
                \node[anchor=east,operator] (app1) at ($(app2.220)+(-1,-0.4)$) {$\App$};
                \node[anchor=east,operator] (nil) at ($(app2.140)+(-2.2,0.8)$) {$\nil$};
                
                \node (in1) at ($(app1)+(-1.3,-0.4)$) {};
                \node (in2) at ($(app1)+(-1.3,0.4)$) {};
                \node (out) at ($(app2.east)+(1,0)$) {};

                \draw[blue] (nil.east) to[out=0,in=140] (app2.140);
                \draw[blue] (in1.east) to[out=0,in=220] (app1.220);
                \draw[blue] (in2.east) to[out=0,in=140] (app1.140);
                \draw[blue] (app1.east) to[out=0,in=220] (app2.220);
                \draw[blue] (app2.east) -- (out) ;
                
                \begin{scope}[on background layer]
                    \node[
                        fit=(nil)(app2),
                        diff,
                        rotate=-20,
                        inner sep=-6pt,
                    ] {};
                \end{scope}
            \end{scope}

            \node at (2,0) {$=$};

            \begin{scope}[xshift=5.5cm]
                \node[operator] (app2) at (0,0) {$\id$};
                \node[anchor=east,operator] (app1) at ($(app2.220)+(-1,-0.4)$) {$\App$};

                \node (in1) at ($(app1)+(-1.3,-0.4)$) {};
                \node (in2) at ($(app1)+(-1.3,0.4)$) {};
                \node (out) at ($(app2.east)+(1,0)$) {};

                \draw[blue] (in1.east) to[out=0,in=220] (app1.220);
                \draw[blue] (in2.east) to[out=0,in=140] (app1.140);
                \draw[blue] (app1.east) to[out=0,in=220] (app2.220);
                \draw[blue] (app2.east) -- (out);
                
                \begin{scope}[on background layer]
                    \node[
                        fit=(app2),
                        result,
                        inner sep=1.5pt,
                    ] {};
                \end{scope}
            \end{scope}

            \node at (7.5,0) {$=$};

            \begin{scope}[xshift=9.5cm]
                \node[operator] (op) at (0,0) {$\App$};

                \node (in1) at ($(op.140)+(-1,0.4)$) {};
                \node (in2) at ($(op.220)+(-1,-0.4)$) {};
                \node (out) at ($(op.east)+(1,0)$) {};

                \draw[blue] (in1.east) to[out=0,in=140] (op.140);
                \draw[blue] (in2.east) to[out=0,in=220] (op.220);
                \draw[blue] (op.east) -- (out);
            \end{scope}
        \end{tikzpicture}
    \end{equation*}
    and for $g$:
    \begin{equation*}
        \begin{tikzpicture}
            \begin{scope}
                \node[operator,fill=Appbg] (app1) at (0,0) {$\App$};
                \node[operator,fill=Appbg] (app2) at ($(app1.140)+(-1.2,0.4)$) {$\App$};
                \node[operator] (nil) at ($(app2.140)+(-1,0.2)$) {$\nil$};

                \node (in1) at ($(app2.220)+(-1,-0.2)$) {};
                \node (in2) at ($(app1.220)+(-2.5,-0.6)$) {};
                \node (out) at ($(app1.east)+(1,0)$) {};

                \draw[blue] (nil.east) to[out=0,in=140] (app2.140);
                \draw[blue] (in1.east) to[out=0,in=220] (app2.220);
                \draw[blue] (in2) to[out=0,in=220] (app1.220);
                \draw[blue] (app1.east) -- (out);
                \draw[blue] (app2.east) to[out=0,in=140] (app1.140);

                \begin{scope}[on background layer]
                    \node[
                        fit=(nil)(app2),
                        diff,
                        rotate=-20,
                        inner sep=-1.5pt,
                    ] {};
                \end{scope}
            \end{scope}

            \node at (2,0) {$=$};

            \begin{scope}[xshift=5.5cm]
                
                \node[operator,fill=Appbg] (app1) at (0,0) {$\App$};
                \node[operator] (app2) at ($(app1.140)+(-1.2,0.4)$) {$\id$};

                \node (in1) at ($(app2.220)+(-1,-0.2)$) {};
                \node (in2) at ($(app1.220)+(-2.5,-0.6)$) {};
                \node (out) at ($(app1.east)+(1,0)$) {};

                \draw[blue] (in1.east) to[out=0,in=220] (app2.220);
                \draw[blue] (in2) to[out=0,in=220] (app1.220);
                \draw[blue] (app1.east) -- (out) ;
                \draw[blue] (app2.east) to[out=0,in=140] (app1.140);

                \begin{scope}[on background layer]
                    \node[
                        fit=(app2),
                        result,
                        inner sep=1.5pt,
                    ] {};
                \end{scope}
            \end{scope}

            \node at (7.5,0) {$=$};

            \begin{scope}[xshift=9.5cm]
                \node[operator] (app) at (0,0) {$\App$};

                \node (in1) at ($(app.140)+(-1,0.4)$) {};
                \node (in2) at ($(app.220)+(-1,-0.4)$) {};
                \node (out) at ($(app.east)+(1,0)$) {};

                \draw[blue] (in1.east) to[out=0,in=140] (op.140);
                \draw[blue] (in2.east) to[out=0,in=220] (op.220);
                \draw[blue] (app.east) -- (out);
            \end{scope}
        \end{tikzpicture}
    \end{equation*}
    Similarly, we have $\cons$-$\beta$ since $\App$ also computes on $\cons$. For $f$:
    \begin{equation*}
        \begin{tikzpicture}
            \begin{scope}
                \node[operator,fill=Appbg] (op) at (0,0) {$\App$};
                \node[operator,fill=consbg] (op1) at ($(op.140)+(-1.5,0.6)$) {$\cons$};
                \node[operator,fill=Appbg] (op2) at ($(op.220)+(-1.5,-0.6)$) {$\App$};
            
                \node (in11) at ($(op1)+(-1.3,0.4)$) {};
                \node (in12) at ($(op1)+(-1.3,-0.4)$) {};
                \node (in21) at ($(op2)+(-1.3,0.4)$) {};
                \node (in22) at ($(op2)+(-1.3,-0.4)$) {};
                \node (out) at ($(op.east)+(1,0)$) {};

                % \draw[blue] (op3.east) to[out=0,in=220] node[inputwire,pos=0.5,yshift=-2pt,below,sloped=false] {$L(A)$} (op1.220);

                \draw[magenta] (in11) to[out=0,in=140] node[inputwire,sloped=false] {} (op1.140);
                \draw[blue] (in12) to[out=0,in=220] node[inputwire,below,sloped=false] {} (op1.220);
                \draw[blue] (in21) to[out=0,in=140] node[inputwire,sloped=false] {} (op2.140);
                \draw[blue] (in22) to[out=0,in=220] node[inputwire,below,sloped=false] {} (op2.220);

                \draw[blue] (op1.east) to[out=0,in=140] node[inputwire,below,sloped=false] {} (op.140);
                \draw[blue] (op2.east) to[out=0,in=220] node[inputwire,below,sloped=false] {} (op.220);

                \draw[blue] (op.east) -- (out) node[outputwire] {};
                \begin{scope}[on background layer]
                    \node[
                        fit=(op1)(op),
                        diff,
                        rotate=-20,
                        inner sep=-3pt,
                    ] {};
                \end{scope}
            \end{scope}
            
            \node at (2,0) {$=$};
            
            % \begin{scope}[xshift=5.5cm]
            %     \node[operator,fill=Appbg] (op) at (0,0) {$\App$};
            %     \node[operator,fill=consbg] (op1) at ($(op.40)+(1.2,0.6)$) {$\cons$};
            %     \node[operator,fill=Appbg] (op2) at ($(op.220)+(-1.2,-0.6)$) {$\App$};

            %     \node (in21) at ($(op2)+(-1.3,0.4)$) {};
            %     \node (in22) at ($(op2)+(-1.3,-0.4)$) {};
            %     \node (out) at ($(op1.east)+(1,0)$) {};

            %     \draw[magenta] (in21 |- 0,1.2) to[out=0,in=140,in looseness=0.2] node[inputwire,sloped=false] {} (op1.140);
            %     \draw[blue] (in21 |- 0,0.4) to[out=0,in=140,in looseness=0.4] node[inputwire,below,sloped=false] {} (op.140);
            %     \draw[blue] (op.east) to[out=0,in=220] node[inputwire,below,sloped=false] {} (op1.220);
            %     \draw[blue] (in21 |- in21) to[out=0,in=140] node[inputwire,sloped=false] {} (op2.140);
            %     \draw[blue] (in22 |- in22) to[out=0,in=220] node[inputwire,below,sloped=false] {} (op2.220);

            %     \draw[blue] (op2.east) to[out=0,in=220] node[inputwire,below,sloped=false] {} (op.220);

            %     \draw[blue] (op1.east) -- (out.west) node[outputwire] {};
            %     \begin{scope}[on background layer]
            %         \node[
            %             fit=(op1)(op),
            %             result,
            %             rotate=20,
            %             inner sep=-3pt,
            %         ] {};
            %         % \node[
            %         %     fit=(op2)(op),
            %         %     rectangle,
            %         %     rounded corners,
            %         %     draw=black!50,
            %         %     inner sep=6pt,
            %         % ] {};
            %         % \node at (0.2,-1) {$f$};
            %     \end{scope}
            % \end{scope}
            \begin{scope}[xshift=5.5cm,yshift=-0.3cm]
                \node[operator,fill=Appbg] (op) at (0,0) {$\App$};
                \node[operator,fill=consbg] (op1) at ($(op.40)+(1.2,0.3)$) {$\cons$};
                \node[operator,fill=Appbg] (op2) at ($(op.220)+(-1.2,-0.3)$) {$\App$};

                \node (in21) at ($(op2)+(-1.3,0.4)$) {};
                \node (in22) at ($(op2)+(-1.3,-0.4)$) {};
                \node (out) at ($(op1.east)+(1,0)$) {};

                \draw[magenta] (in21 |- 0,1.5) to[out=0,in=140] node[inputwire,sloped=false] {} (op1.140);
                \draw[blue] (in21 |- 0,0.7) to[out=0,in=140] node[inputwire,below,sloped=false] {} (op.140);
                \draw[blue] (op.east) to[out=0,in=220] node[inputwire,below,sloped=false] {} (op1.220);
                \draw[blue] (in21 |- in21) to[out=0,in=140] node[inputwire,sloped=false] {} (op2.140);
                \draw[blue] (in22 |- in22) to[out=0,in=220] node[inputwire,below,sloped=false] {} (op2.220);

                \draw[blue] (op2.east) to[out=0,in=220] node[inputwire,below,sloped=false] {} (op.220);

                \draw[blue] (op1.east) -- (out.west) node[outputwire] {};
                \begin{scope}[on background layer]
                    \node[
                        fit=(op1)(op),
                        result,
                        rotate=15,
                        inner sep=-3pt,
                    ] {};
                    % \node[
                    %     fit=(op2)(op),
                    %     rectangle,
                    %     rounded corners,
                    %     draw=black!50,
                    %     inner sep=6pt,
                    % ] {};
                    % \node at (0.2,-1) {$f$};
                \end{scope}
            \end{scope}
        \end{tikzpicture}
    \end{equation*}
    and for $g$, we push the $\cons$ through $\App$ with two computations:
    \begin{equation*}
        \begin{tikzpicture}
            \begin{scope}[yshift=0.3cm]
                \node[operator,fill=Appbg] (op) at (0,0) {$\App$};
                \node[operator,fill=consbg] (op1) at ($(op.140)+(-1.2,0.3)$) {$\cons$};
                \node[operator,fill=Appbg] (op2) at ($(op.340)+(1.2,-0.3)$) {$\App$};

                \node (in1) at ($(op1.140)+(-1,0.2)$) {};
                \node (in2) at ($(op1.220)+(-1,-0.2)$) {};
                \node (out) at ($(op2.east)+(1,0)$) {};

                \draw[magenta] (in1 |- in1) to[out=0,in=140] node[inputwire] {} (op1.140);
                \draw[blue] (in2 |- in2) to[out=0,in=220] node[inputwire] {} (op1.220);
                \draw[blue] (in2 |- 0,-0.7) to[out=0,in=220] node[inputwire] {} (op.220);
                \draw[blue] (in2 |- 0,-1.5) to[out=0,in=220] node[inputwire] {} (op2.220);
                
                \draw[blue] (op1.east) to[out=0,in=140] node[inputwire] {} (op.140);
                \draw[blue] (op.east) to[out=0,in=140] node[inputwire] {} (op2.140);
                \draw[blue] (op2.east) -- (out.west) node[outputwire] {};
                \begin{scope}[on background layer]
                    \node[
                        fit=(op1)(op),
                        diff,
                        rotate=-15,
                        inner sep=-3pt,
                    ] {};
                \end{scope}
            \end{scope}
                
            \node at (3,0) {$=$};
            
            \begin{scope}[xshift=6.2cm,yshift=0.5cm]
                \node[operator,fill=consbg] (op) at (0,0) {$\cons$};
                \node[operator,fill=Appbg] (op1) at ($(op.220)+(-1.2,-0.3)$) {$\App$};
                \node[operator,fill=Appbg] (op2) at ($(op.340)+(1.2,-0.3)$) {$\App$};

                \node (in1) at ($(op1.140)+(-1,0.2)$) {};
                \node (in2) at ($(op1.220)+(-1,-0.2)$) {};
                \node (out) at ($(op2.east)+(1,0)$) {};

                \draw[magenta] (in2 |- 0,0.7) to[out=0,in=140] node[inputwire] {} (op.140);
                \draw[blue] (in1 |- in1) to[out=0,in=140] node[inputwire] {} (op1.140);
                \draw[blue] (in2 |- in2) to[out=0,in=220] node[inputwire] {} (op1.220);
                \draw[blue] (in2 |- 0,-1.7) to[out=0,in=220] node[inputwire] {} (op2.220);
                
                \draw[blue] (op1.east) to[out=0,in=220] node[inputwire] {} (op.220);
                \draw[blue] (op.east) to[out=0,in=140] node[inputwire] {} (op2.140);
                \draw[blue] (op2.east) -- (out.west) node[outputwire] {};
                \begin{scope}[on background layer]
                    \node[
                        fit=(op1)(op),
                        result,
                        rotate=15,
                        inner sep=-3pt,
                    ] {};
                    \node[
                        fit=(op)(op2),
                        diff,
                        rotate=-15,
                        inner sep=-3pt,
                    ] {};
                \end{scope}
            \end{scope}

            \node at (9.4,0) {$=$};

            \begin{scope}[xshift=12.6cm,yshift=-0.5cm]
                \node[operator,fill=Appbg] (op) at (0,0) {$\App$};
                \node[operator,fill=Appbg] (op1) at ($(op.140)+(-1.2,0.3)$) {$\App$};
                \node[operator,fill=consbg] (op2) at ($(op.40)+(1.2,0.3)$) {$\cons$};

                \node (in1) at ($(op1.140)+(-1,0.2)$) {};
                \node (in2) at ($(op1.220)+(-1,-0.2)$) {};
                \node (out) at ($(op2.east)+(1,0)$) {};

                \draw[magenta] (in2 |- 0,1.7) to[out=0,in=140] node[inputwire] {} (op2.140);
                \draw[blue] (in1 |- in1) to[out=0,in=140] node[inputwire] {} (op1.140);
                \draw[blue] (in2 |- in2) to[out=0,in=220] node[inputwire] {} (op1.220);
                \draw[blue] (in2 |- 0,-0.7) to[out=0,in=220] node[inputwire] {} (op.220);
                
                \draw[blue] (op1.east) to[out=0,in=140] node[inputwire] {} (op.140);
                \draw[blue] (op.east) to[out=0,in=220] node[inputwire] {} (op2.220);
                \draw[blue] (op2.east) -- (out.west) node[outputwire] {};
                \begin{scope}[on background layer]
                    \node[
                        fit=(op)(op2),
                        result,
                        rotate=20,
                        inner sep=-3pt,
                    ] {};
                \end{scope}
            \end{scope}
        \end{tikzpicture}
    \end{equation*}
    Hence $f=\iter(\App,\cons)=g$.

    It remains to show the unique extension property. Let $f:A\to B$ be a morphism, with $(B,e,m)$ a monoid object. We define the extension $f^\sharp:L(A)\to B$ as the mediating map $\iter(n,c)$, where $n:1\rightarrow B$ and $c:A\otimes B\rightarrow B$ are:
    \begin{equation*}
        \begin{tikzpicture}
            \begin{scope}
                \node[operator] (op) at (-1,0) {$e$};
                \node (out) at ($(op.east)+(1,0)$) {};
                \draw[green] (op.east) -- node[outputwire] {$B$} (out);
            \end{scope}
            \begin{scope}[xshift=6cm]
                \node[operator] (m) at (0,0) {$m$};
                \node[operator] (f) at ($(m.140)+(-1.5,0.4)$) {$f$};
                \node (in1) at ($(f.west)+(-1,0)$) {};
                \node (in2) at ($(in1 |- m.220) + (0,-0.4)$) {};
                \node (out) at ($(m.east)+(1,0)$) {};

                \draw[magenta] (in1) -- node[inputwire] {$A$} (f);
                \draw[green] (f) to[out=0,in=140] node[above] {$B$} (m.140);
                \draw[green] (in2) to[out=0,in=220,in looseness=0.6] node[inputwire,pos=0.1,below] {$B$} (m.220);
                \draw[green] (m) -- node[inputwire] {$B$} (out);
            \end{scope}
        \end{tikzpicture}
    \end{equation*}
    and the unit map $\eta:A\to L(A)$ as:
    \begin{equation*}
        \begin{tikzpicture}
            \node[operator,fill=consbg] (op) at (0,0) {$\cons$};
            \node (in1) at ($(op.140)+(-1.5,0.4)$) {};
            \node[anchor=east,operator] (in2) at ($(op.220)+(-1,-0.4)$) {$\nil$};
            \node (out) at ($(op.east)+(1,0)$) {};
            \draw[magenta] (in1) to[out=0,in=140] node[inputwire,sloped=false] {$A$} (op.140);
            \draw[blue] (in2.east) to[out=0,in=220] node[inputwire,below,pos=0.4,sloped=false] {$L(A)$} (op.220);
            \draw[blue] (op.east) -- (out) node[outputwire] {$L(A)$};
        \end{tikzpicture}
    \end{equation*}
    Note that $\App$ computes on $\eta$ in the left argument as:
    \begin{equation*}
        \begin{tikzpicture}
            \begin{scope}
                \node[operator] (app) at (0,0) {$\App$};
                \node[operator] (eta) at ($(app.140)+(-1.2,0.4)$) {$\eta$};
                \node (in1) at ($(eta.west)+(-1,0)$) {};
                \node (in2) at ($(in1 |- app.220) + (0,-0.4)$) {};
                \node (out) at ($(app.east)+(1,0)$) {};

                \draw[magenta] (in1) -- (eta);
                \draw[blue] (eta) to[out=0,in=140] (app.140);
                \draw[blue] (in2) to[out=0,in=220,in looseness=0.6] (app.220);
                \draw[blue] (app) -- (out);
            \end{scope}
            \node at (2,0) {$=$};
            \begin{scope}[xshift=4.2cm]
                \node[operator,fill=consbg] (op) at (0,0) {$\cons$};
                \node (in1) at ($(op.140)+(-1.2,0.4)$) {};
                \node (in2) at ($(op.220)+(-1.2,-0.4)$) {};
                \node (out) at ($(op.east)+(1,0)$) {};
                \draw[magenta] (in1) to[out=0,in=140] (op.140);
                \draw[blue] (in2.east) to[out=0,in=220] (op.220);
                \draw[blue] (op.east) -- (out);
            \end{scope}
        \end{tikzpicture}
    \end{equation*}
    since $\App$-$\cons$-$\beta$ commutes the $\App$ next to the $\nil$, from which it reduces to the identity by $\App$-$\nil$-$\beta$. (See Lemma \ref{lem:eta-app}.)

    We verify that $f^\sharp\circ\eta=f$:
    \begin{equation*}
        \begin{tikzpicture}
            \begin{scope}
                \node[operator,inner sep=1pt] (fs) at (0,0) {$f^\sharp$};
                \node[operator] (eta) at (-1.5,0) {$\eta$};

                \node (in) at ($(eta.west)+(-1,0)$) {};
                \node (out) at ($(fs.east)+(1,0)$) {};

                \draw[magenta] (in) -- node[inputwire] {} (eta);
                \draw[blue] (eta) -- node[inputwire] {} (fs);
                \draw[green] (fs) -- node[inputwire] {} (out);
            \end{scope}
            
            \node at (2,0) {$=$};
            
            \begin{scope}[xshift=6cm]
                \node[operator,inner sep=1pt] (fs) at (0,0) {$f^\sharp$};
                \node[operator] (cons) at (-1.5,0) {$\cons$};

                \node (in) at ($(cons.140)+(-1.3,0.4)$) {};
                \node[operator] (nil) at ($(cons.220)+(-1,-0.4)$) {$\nil$};
                \node (out) at ($(fs.east)+(1,0)$) {};

                \draw[magenta] (in) to[out=0,in=140] node[inputwire] {} (cons.140);
                \draw[blue] (nil) to[out=0,in=220] node[inputwire] {} (cons.220);
                \draw[blue] (cons) -- node[inputwire] {} (fs);
                \draw[green] (fs) -- node[inputwire] {} (out);

                \node[anchor=west] at ($(out)+(0.2,0)$) {$f^\sharp$-$\cons$-$\beta$};

                \begin{scope}[on background layer]
                    \node[
                        fit=(cons)(fs),
                        diff,
                        minimum height=0.9cm,
                        inner sep=-4.5pt,
                    ] {};
                \end{scope}
            \end{scope}
            
            \node at (2,-2.5) {$=$};
            
            \begin{scope}[xshift=6cm,yshift=-2.5cm]
                \node[operator] (m) at (0,0) {$m$};
                \node[operator,inner sep=1pt] (fs) at ($(m.220)+(-1.3,-0.4)$) {$f^\sharp$};
                \node[operator] (f) at ($(m.140)+(-1.3,0.4)$) {$f$};
                \node[operator] (nil) at (nil |- fs) {$\nil$};

                \node (in) at (in |- f) {};
                \node (out) at (out |- m) {};

                \draw[magenta] (in) -- (f);
                \draw[green] (f) to[out=0,in=140] (m.140);
                \draw[blue] (nil) -- (fs);
                \draw[green] (fs) to[out=0,in=220] (m.220);
                \draw[green] (m) -- (out);

                \node[anchor=west] at ($(out)+(0.2,0)$) {$f^\sharp$-$\nil$-$\beta$};

                \begin{scope}[on background layer]
                    \node[
                        fit=(nil)(fs),
                        diff,
                        minimum height=0.9cm,
                        inner sep=-4.5pt,
                    ] {};
                \end{scope}
            \end{scope}

            \node at (2,-5) {$=$};

            \begin{scope}[xshift=6cm,yshift=-5cm]
                \node[operator] (m) at (0,0) {$m$};
                \node[operator,inner sep=1pt] (e) at ($(m.220)+(-1.3,-0.4)$) {$e$};
                \node[operator] (f) at ($(m.140)+(-1.3,0.4)$) {$f$};

                \node (in) at (in |- f) {};
                \node (out) at (out |- m) {};

                \draw[magenta] (in) -- (f);
                \draw[green] (f) to[out=0,in=140] (m.140);
                \draw[green] (e) to[out=0,in=220] (m.220);
                \draw[green] (m) -- (out);

                \node[anchor=west] at ($(out)+(0.2,0)$) {$B$-right unit};

                \begin{scope}[on background layer]
                    \node[
                        fit=(e),
                        result,
                        inner sep=1.5pt,
                    ] {};
                    \node[
                        fit=(e)(m),
                        diff,
                        rotate=20,
                        inner sep=-1.5pt,
                    ] {};
                \end{scope}
            \end{scope}

            \node at (2,-7.5) {$=$};

            \begin{scope}[xshift=6cm,yshift=-7.75cm]
                \node[operator] (id) at (0,0) {$\id$};
                \node[operator] (f) at ($(id.140)+(-1.3,0.4)$) {$f$};

                \node (in) at (in |- f) {};
                \node (out) at (out |- id) {};

                \draw[magenta] (in) -- (f);
                \draw[green] (f) to[out=0,in=140] (id.140);
                \draw[green] (id) -- (out);
                \begin{scope}[on background layer]
                    \node[
                        fit=(id),
                        result,
                        inner sep=1.5pt,
                    ] {};
                \end{scope}
            \end{scope}

            \node at (2,-9) {$=$};

            \begin{scope}[xshift=6cm,yshift=-9cm]
                \node[operator] (f) at (-1.55,0) {$f$};

                \node (in) at (in |- f) {};
                \node (out) at (out |- f) {};

                \draw[magenta] (in) -- (f);
                \draw[green] (f) -- (out);
            \end{scope}
        \end{tikzpicture}
    \end{equation*}
    Next, we check that $f^\sharp$ is a monoid homomorphism.
    
    Since $f^\sharp=\iter(e,\ldots)$, $f^\sharp$ computes on $\nil$ to $e$ and hence preserves the monoid identity. For the homomorphism property, let $g,h:L(A)\otimes L(A)\to B$ be the following maps:
    \begin{equation*}
        \begin{tikzpicture}
            \begin{scope}
                \node[operator,inner sep=1pt] (fs) at (0,0) {$f^\sharp$};
                \node[operator] (app) at (-1.5,0) {$\App$};

                \node (in1) at ($(app.140)+(-1.3,0.4)$) {};
                \node (in2) at ($(app.220)+(-1.3,-0.4)$) {};
                \node (out) at ($(fs.east)+(1,0)$) {};

                \draw[blue] (in1) to[out=0,in=140] node[inputwire] {$L(A)$} (app.140);
                \draw[blue] (in2) to[out=0,in=220] node[inputwire,below] {$L(A)$} (app.220);
                \draw[blue] (app) -- node[above] {$L(A)$} (fs);
                \draw[green] (fs) -- node[outputwire] {$B$} (out);
            \end{scope}
            \begin{scope}[xshift=6cm]
                \node[operator] (m) at (0,0) {$m$};
                \node[operator,inner sep=1pt] (fs1) at ($(m.140)+(-1.3,0.4)$) {$f^\sharp$};
                \node[operator,inner sep=1pt] (fs2) at ($(m.220)+(-1.3,-0.4)$) {$f^\sharp$};

                \node (in1) at ($(fs1.west)+(-1,0)$) {};
                \node (in2) at ($(fs2.west)+(-1,0)$) {};
                \node (out) at ($(m.east)+(1,0)$) {};

                \draw[blue] (in1) -- node[inputwire] {$L(A)$} (fs1);
                \draw[blue] (in2) -- node[inputwire,below] {$L(A)$} (fs2);
                \draw[green] (fs1) to[out=0,in=140] node[inputwire] {$B$} (m.140);
                \draw[green] (fs2) to[out=0,in=220] node[inputwire,below] {$B$} (m.220);
                \draw[green] (m) -- node[outputwire] {$B$} (out);
            \end{scope}
        \end{tikzpicture}
    \end{equation*}
    We show that these maps are both the iterator $\iter(f^\sharp,c)$, where $c$ is the same list algebra map as in the definition of $f^\sharp$ above. We have $g$-$\nil$-$\beta$ by $\App$-$\nil$-$\beta$, and $h$-$\nil$-$\beta$ by $f^\sharp$-$\nil$-$\beta$ and left unit for $B$:
    \begin{equation*}
        \begin{tikzpicture}[baseline=(g)]
            \begin{scope}
                \node[operator] (g) at (0,0) {$g$};
                \node[anchor=east,operator] (nil) at ($(g.140)+(-1,0.4)$) {$\nil$};
                \node (in) at ($(g.220)+(-1.5,-0.4)$) {};
                \node (out) at ($(g.east)+(1,0)$) {};
                \draw[blue] (nil.east) to[out=0,in=140] (g.140);
                \draw[blue] (in) to[out=0,in=220] (g.220);
                \draw[green] (g.east) -- (out);
            \end{scope}

            \node at (2,0) {$=$};

            \begin{scope}[xshift=6cm]
                \node[operator,inner sep=1pt] (fs) at (0,0) {$f^\sharp$};
                \node[operator] (app) at (-1.5,0) {$\App$};

                \node[anchor=east,operator] (nil) at ($(app.140)+(-1,0.4)$) {$\nil$};
                \node (in) at ($(app.220)+(-1.5,-0.4)$) {};
                \node (out) at ($(fs.east)+(1,0)$) {};

                \draw[blue] (nil) to[out=0,in=140] (app.140);
                \draw[blue] (in) to[out=0,in=220] (app.220);
                \draw[blue] (app) -- (fs);
                \draw[green] (fs) -- (out);
                
                \begin{scope}[on background layer]
                    \node[
                        fit=(nil)(app),
                        diff,
                        rotate=-20,
                        inner sep=-3pt,
                    ] {};
                \end{scope}
            \end{scope}

            \node at (2,-2) {$=$};

            \begin{scope}[xshift=6cm,yshift=-2cm]
                \node[operator,inner sep=1pt] (fs) at (0,0) {$f^\sharp$};
                \node[operator] (id) at (-1.5,0) {$\id$};

                \node (in) at ($(in |- id.220)+(0,-0.4)$) {};
                \node (out) at (out |- fs) {};

                \draw[blue] (in) to[out=0,in=220] (id.220);
                \draw[blue] (id) -- (fs);
                \draw[green] (fs) -- (out);

                \begin{scope}[on background layer]
                    \node[
                        fit=(id),
                        result,
                        inner sep=1.5pt,
                    ] {};
                \end{scope}
            \end{scope}

            \node at (2,-4) {$=$};

            \begin{scope}[xshift=6cm,yshift=-4cm]
                \node[operator,inner sep=1pt] (fs) at (0,0) {$f^\sharp$};

                \node (in) at (in |- fs) {};
                \node (out) at (out |- fs) {};

                \draw[blue] (in) -- (fs);
                \draw[green] (fs) -- (out);
            \end{scope}
        \end{tikzpicture}
        \qquad % ===============================================================
        \begin{tikzpicture}[baseline=(h)]
            \begin{scope}
                \node[operator] (h) at (0,0) {$h$};
                \node[anchor=east,operator] (nil) at ($(h.140)+(-1,0.4)$) {$\nil$};
                \node (in) at ($(h.220)+(-1.5,-0.4)$) {};
                \node (out) at ($(h.east)+(1,0)$) {};
                \draw[blue] (nil.east) to[out=0,in=140] (h.140);
                \draw[blue] (in) to[out=0,in=220] (h.220);
                \draw[green] (h.east) -- (out);
            \end{scope}

            \node at (2,0) {$=$};

            \begin{scope}[xshift=6cm]
                \node[operator] (m) at (0,0) {$m$};
                \node[operator,inner sep=1pt] (fs1) at ($(m.140)+(-1.3,0.4)$) {$f^\sharp$};
                \node[operator,inner sep=1pt] (fs2) at ($(m.220)+(-1.3,-0.4)$) {$f^\sharp$};
                \node[operator] (nil) at ($(fs1.west)+(-1,0)$) {$\nil$};

                \node (in) at ($(fs2.west)+(-1.3,0)$) {};
                \node (out) at ($(m.east)+(1,0)$) {};

                \draw[blue] (nil) -- (fs1);
                \draw[blue] (in) -- (fs2);
                \draw[green] (fs1) to[out=0,in=140] (m.140);
                \draw[green] (fs2) to[out=0,in=220] (m.220);
                \draw[green] (m) -- (out);

                \begin{scope}[on background layer]
                    \node[
                        fit=(fs1)(nil),
                        diff,
                        minimum height=0.9cm,
                        inner sep=-4.5pt,
                    ] {};
                \end{scope}
            \end{scope}

            \node at (2,-2.5) {$=$};

            \begin{scope}[xshift=6cm,yshift=-2.5cm]
                \node[operator] (m) at (0,0) {$m$};
                \node[operator,inner sep=1pt] (e) at ($(m.140)+(-1.3,0.4)$) {$e$};
                \node[operator,inner sep=1pt] (fs) at ($(m.220)+(-1.3,-0.4)$) {$f^\sharp$};

                \node (in) at (in |- fs) {};
                \node (out) at (out |- m) {};

                \draw[blue] (in) -- (fs);
                \draw[green] (e) to[out=0,in=140] (m.140);
                \draw[green] (fs) to[out=0,in=220] (m.220);
                \draw[green] (m) -- (out);

                \begin{scope}[on background layer]
                    \node[
                        fit=(e),
                        result,
                        inner sep=1.5pt,
                    ] {};
                    \node[
                        fit=(e)(m),
                        diff,
                        rotate=-20,
                        inner sep=-1.5pt,
                    ] {};
                \end{scope}
            \end{scope}

            \node at (2,-4.5) {$=$};

            \begin{scope}[xshift=6cm,yshift=-4.25cm]
                \node[operator] (id) at (0,0) {$id$};
                \node[operator,inner sep=1pt] (fs) at ($(id.220)+(-1.3,-0.4)$) {$f^\sharp$};

                \node (in) at (in |- fs) {};
                \node (out) at (out |- id) {};

                \draw[blue] (in) -- (fs);
                \draw[green] (fs) to[out=0,in=220] (id.220);
                \draw[green] (id) -- (out);

                \begin{scope}[on background layer]
                    \node[
                        fit=(id),
                        result,
                        inner sep=1.5pt,
                    ] {};
                \end{scope}
            \end{scope}

            \node at (2,-5.75) {$=$};

            \begin{scope}[xshift=6cm,yshift=-5.75cm]
                \node[operator,inner sep=1pt] (fs) at (fs |- 0,0) {$f^\sharp$};

                \node (in) at (in |- fs) {};
                \node (out) at (out |- fs) {};

                \draw[blue] (in) -- (fs);
                \draw[green] (fs) -- (out);
            \end{scope}
        \end{tikzpicture}
    \end{equation*}
    For $\cons$-$\beta$ we have the computations:
    \begin{equation*}
        \begin{tikzpicture}[baseline=(g),scale=0.8]
            \begin{scope}
                \node[operator] (g) at (0,0) {$g$};
                \node[anchor=east,operator,fill=consbg] (cons) at ($(g.140)+(-1,0.4)$) {$\cons$};
                
                \node (in1) at ($(cons.140)+(-1,0.2)$) {};
                \node (in2) at ($(cons.220)+(-1,-0.2)$) {};
                \node (in3) at ($(in1 |- g.220)+(0,-0.6)$) {};
                \node (out) at ($(g.east)+(1,0)$) {};

                \draw[magenta] (in1) to[out=0,in=140] (cons.140);
                \draw[blue] (in2) to[out=0,in=220] (cons.220);
                \draw[blue] (in3) to[out=0,in=220] (g.220);
                \draw[blue] (cons) to[out=0,in=140] (g.140);
                \draw[green] (g) -- (out);
            \end{scope}

            \node at (1.7,0) {$=$};

            \begin{scope}[xshift=5.1cm]
                \node[operator,fill=Appbg] (app) at (0,0) {$\App$};
                \node[anchor=east,operator,fill=Appbg] (cons) at ($(app.140)+(-1,0.4)$) {$\cons$};
                \node[operator,inner sep=1pt] (fs) at ($(app)+(1.5,0)$) {$f^\sharp$};
                
                \node (in1) at ($(cons.140)+(-1,0.2)$) {};
                \node (in2) at ($(cons.220)+(-1,-0.2)$) {};
                \node (in3) at ($(in1 |- app.220)+(0,-0.6)$) {};
                \node (out) at ($(fs.east)+(1,0)$) {};

                \draw[magenta] (in1) to[out=0,in=140] (cons.140);
                \draw[blue] (in2) to[out=0,in=220] (cons.220);
                \draw[blue] (in3) to[out=0,in=220] (app.220);
                \draw[blue] (cons) to[out=0,in=140] (app.140);
                \draw[blue] (app) -- (fs);
                \draw[green] (fs) -- (out);

                \begin{scope}[on background layer]
                    \node[
                        fit=(cons)(app),
                        diff,
                        rotate=-20,
                        inner sep=-3pt,
                    ] {};
                \end{scope}
            \end{scope}

            \node at (1.7,-2.5) {$=$};

            \begin{scope}[xshift=5.1cm,yshift=-2.5cm]
                \node[operator,fill=consbg] (cons) at (0,0) {$\cons$};
                \node[anchor=east,operator,fill=consbg] (app) at ($(cons.220)+(-1,-0.4)$) {$\App$};
                \node[operator,inner sep=1pt] (fs) at ($(cons)+(1.5,0)$) {$f^\sharp$};

                \node (in1) at ($(in1 |- cons.140)+(0,0.6)$) {};
                \node (in2) at ($(in1 |- app.140)+(0,0.2)$) {};
                \node (in3) at ($(in1 |- app.220)+(0,-0.2)$) {};
                \node (out) at (out |- fs) {};

                \draw[magenta] (in1) to[out=0,in=140] (cons.140);
                \draw[blue] (in2) to[out=0,in=140] (app.140);
                \draw[blue] (in3) to[out=0,in=220] (app.220);
                \draw[blue] (app) to[out=0,in=220] (cons.220);
                \draw[blue] (cons) -- (fs);
                \draw[green] (fs) -- (out);

                \begin{scope}[on background layer]
                    \node[
                        fit=(cons)(app),
                        result,
                        rotate=20,
                        inner sep=-3pt,
                    ] {};
                    \node[
                        fit=(cons)(fs),
                        diff,
                        minimum height=0.8cm,
                        inner sep=-3pt,
                    ] {};
                \end{scope}
            \end{scope}

            \node at (1.7,-5) {$=$};

            \begin{scope}[xshift=5.1cm,yshift=-5cm]
                \node[operator] (m) at (fs |- 0,0) {$m$};
                \node[operator,inner sep=1pt] (fs) at ($(cons |- m.220)+(0,-0.4)$) {$f^\sharp$};
                \node[operator] (f) at ($(cons |- m.140)+(0,0.4)$) {$f$};
                \node[operator,fill=consbg] (app) at (app |- fs) {$\App$};

                \node (in1) at (in1 |- f) {};
                \node (in2) at ($(in2 |- app.140)+(0,0.2)$) {};
                \node (in3) at ($(in3 |- app.220)+(0,-0.2)$) {};
                \node (out) at (out |- m) {};

                \draw[magenta] (in1) -- (f);
                \draw[blue] (in2) to[out=0,in=140] (app.140);
                \draw[blue] (in3) to[out=0,in=220] (app.220);
                \draw[green] (f) to[out=0,in=140] (m.140);
                \draw[green] (fs) to[out=0,in=220] (m.220);
                \draw[green] (m) -- (out);
                \draw[blue] (app) -- (fs);

                \begin{scope}[on background layer]
                    \node[
                        fit=(f)(m),
                        result,
                        rotate=-25,
                        inner sep=-3pt,
                    ] {};
                \end{scope}
            \end{scope}

            \node at (1.7,-7.5) {$=$};

            \begin{scope}[xshift=6.4cm,yshift=-7.5cm]
                \node[operator] (m) at (m |- 0,0) {$m$};
                \node[operator] (f) at ($(f |- m.140)+(0,0.4)$) {$f$};
                \node[operator] (g) at ($(m.220)+(-2,-0.4)$) {$g$};

                \node (in1) at (in1 |- f) {};
                \node (in2) at ($(in2 |- g.140)+(0,0.2)$) {};
                \node (in3) at ($(in3 |- g.220)+(0,-0.2)$) {};
                \node (out) at (out |- m) {};

                \draw[magenta] (in1) -- (f);
                \draw[blue] (in2) to[out=0,in=140,in looseness=0.7] (g.140);
                \draw[blue] (in3) to[out=0,in=220,in looseness=0.7] (g.220);
                \draw[green] (f) to[out=0,in=140] (m.140);
                \draw[green] (m) -- (out);
                \draw[green] (g) to[out=0,in=220,in looseness=0.8] (m.220);
            \end{scope}
        \end{tikzpicture}
        \qquad % ===============================================================
        \begin{tikzpicture}[baseline=(h),scale=0.8]
            \begin{scope}
                \node[operator] (h) at (0,0) {$h$};
                \node[operator,fill=consbg] (cons) at ($(h.140)+(-1.5,0.4)$) {$\cons$};
                
                \node (in1) at ($(cons.140)+(-1,0.2)$) {};
                \node (in2) at ($(cons.220)+(-1,-0.2)$) {};
                \node (in3) at ($(in1 |- h.220)+(0,-0.6)$) {};
                \node (out) at ($(h.east)+(1,0)$) {};

                \draw[magenta] (in1) to[out=0,in=140] (cons.140);
                \draw[blue] (in2) to[out=0,in=220] (cons.220);
                \draw[blue] (in3) to[out=0,in=220] (h.220);
                \draw[blue] (cons) to[out=0,in=140] (h.140);
                \draw[green] (h) -- (out);
            \end{scope}

            \node at (1.7,0) {$=$};

            \begin{scope}[xshift=6.2cm]
                \node[operator] (m) at (0,0) {$m$};
                \node[operator,inner sep=1pt] (fs1) at ($(m.140)+(-1,0.4)$) {$f^\sharp$};
                \node[operator,fill=consbg] (cons) at ($(fs1) + (-1.5,0)$) {$\cons$};
                \node[operator,inner sep=1pt] (fs2) at ($(cons |- m.220)+(0,-0.4)$) {$f^\sharp$};

                \node (in1) at ($(cons.140)+(-1,0.2)$) {};
                \node (in2) at ($(cons.220)+(-1,-0.2)$) {};
                \node (in3) at (in1 |- fs2) {};
                \node (out) at ($(m.east)+(1,0)$) {};

                \draw[magenta] (in1) to[out=0,in=140] (cons);
                \draw[blue] (in2) to[out=0,in=220] (cons);
                \draw[blue] (in3) -- (fs2);
                \draw[blue] (cons) -- (fs1);
                \draw[green] (fs1) to[out=0,in=140] (m.140);
                \draw[green] (fs2) to[out=0,in=220] (m.220);
                \draw[green] (m) -- (out);

                \begin{scope}[on background layer]
                    \node[
                        fit=(cons)(fs1),
                        diff,
                        minimum height=0.8cm,
                        inner sep=-3pt,
                    ] {};
                \end{scope}
            \end{scope}

            \node at (1.7,-2.5) {$=$};

            \begin{scope}[xshift=6.2cm,yshift=-2.7cm]
                \node[operator] (m2) at (0,0) {$m$};
                \node[operator] (m1) at ($(m2.140)+(-1,0.4)$) {$m$};
                \node[operator,inner sep=1pt] (fs2) at ($(cons |- m2.220)+(0,-0.4)$) {$f^\sharp$};
                \node[operator] (f) at ($(cons |- m1.140)+(0,0.2)$) {$f$};
                \node[operator,inner sep=1pt] (fs1) at ($(cons |- m1.220)+(0,-0.2)$) {$f^\sharp$};
                
                \node (in1) at (in1 |- f) {};
                \node (in2) at (in2 |- fs1) {};
                \node (in3) at (in1 |- fs2) {};
                \node (out) at (out |- m2) {};

                \draw[magenta] (in1) -- (f);
                \draw[blue] (in2) -- (fs1);
                \draw[blue] (in3) -- (fs2);
                \draw[green] (f) to[out=0,in=140] (m1.140);
                \draw[green] (fs1) to[out=0,in=220] (m1.220);
                \draw[green] (m1) to[out=0,in=140] (m2.140);
                \draw[green] (fs2) to[out=0,in=220] (m2.220);
                \draw[green] (m2) -- (out);

                \begin{scope}[on background layer]
                    \coordinate (F1) at ($(f)+(-0.5,0.7)$);
                    \coordinate (F2) at ($(fs1)+(-0.5,-0.7)$);
                    \coordinate (F3) at ($(m1)+(1.3,0)$);
                    \draw[result,rounded corners=12pt] (F1) -- (F2) [rounded corners=25pt] -- (F3) [rounded corners=12pt] -- cycle;
                    \node[
                        fit=(m1)(m2),
                        diff,
                        rotate=-25,
                        inner sep=-2pt,
                    ] {};
                \end{scope}
            \end{scope}

            \node at (1.7,-5) {$=$};

            \begin{scope}[xshift=6.2cm,yshift=-4.8cm]
                \node[operator] (m2) at (0,0) {$m$};
                \node[operator] (m1) at ($(m2.220)+(-1,-0.4)$) {$m$};
                \node[operator] (f) at ($(f |- m2.140)+(0,0.4)$) {$f$};
                \node[operator,inner sep=1pt] (fs1) at ($(fs1 |- m1.140)+(0,0.2)$) {$f^\sharp$};
                \node[operator,inner sep=1pt] (fs2) at ($(fs2 |- m1.220)+(0,-0.2)$) {$f^\sharp$};
                
                \node (in1) at (in1 |- f) {};
                \node (in2) at (in2 |- fs1) {};
                \node (in3) at (in1 |- fs2) {};
                \node (out) at (out |- m2) {};

                \draw[magenta] (in1) -- (f);
                \draw[blue] (in2) -- (fs1);
                \draw[blue] (in3) -- (fs2);
                \draw[green] (fs1) to[out=0,in=140] (m1.140);
                \draw[green] (fs2) to[out=0,in=220] (m1.220);
                \draw[green] (f) to[out=0,in=140] (m2.140);
                \draw[green] (m1) to[out=0,in=220] (m2.220);
                \draw[green] (m2) -- (out);

                \begin{scope}[on background layer]
                    \node[
                        fit=(m1)(m2),
                        diff,
                        rotate=30,
                        inner sep=-3pt,
                    ] {};
                \end{scope}
            \end{scope}

            \node at (1.7,-7.5) {$=$};

            \begin{scope}[xshift=6.2cm,yshift=-7.3cm]
                \node[operator] (m) at (0,0) {$m$};
                \node[operator] (f) at ($(f |- m.140)+(0,0.4)$) {$f$};
                \node[operator] (h) at ($(m.220)+(-1.8,-0.4)$) {$h$};

                \node (in1) at (in1 |- f) {};
                \node (in2) at ($(in2 |- h.140)+(0,0.2)$) {};
                \node (in3) at ($(in3 |- h.220)+(0,-0.2)$) {};
                \node (out) at (out |- m) {};

                \draw[magenta] (in1) -- (f);
                \draw[blue] (in2) to[out=0,in=140,in looseness=0.7] (h.140);
                \draw[blue] (in3) to[out=0,in=220,in looseness=0.7] (h.220);
                \draw[green] (f) to[out=0,in=140,in looseness=0.8] (m.140);
                \draw[green] (m) -- (out);
                \draw[green] (h) to[out=0,in=220,in looseness=0.8] (m.220);
            \end{scope}
        \end{tikzpicture}
    \end{equation*}
    Finally, we show that the extension $f^\sharp$ is unique. Suppose there is another monoid homomorphism $g:L(A)\to B$ satisfying $g\circ\eta=f$. Then, $g$ satisfies the same computation rules as $f^\sharp$. Indeed, $\nil$-$\beta$ holds since $g$ preserves identities, and $\cons$-$\beta$ holds via the following computation:
    \begin{equation*}
        \begin{tikzpicture}
            \begin{scope}
                \node[operator,] (g) at (0,0) {$g$};
                \node[operator] (cons) at (-1.5,0) {$\cons$};

                \node (in1) at ($(cons.140)+(-1.3,0.4)$) {};
                \node (in2) at ($(cons.220)+(-1.3,-0.4)$) {};
                \node (out) at ($(g.east)+(1,0)$) {};

                \draw[magenta] (in1) to[out=0,in=140] (cons.140);
                \draw[blue] (in2) to[out=0,in=220] (cons.220);
                \draw[blue] (cons) --  (g);
                \draw[green] (g) -- (out);

                \begin{scope}[on background layer]
                    \node[
                        fit=(cons),
                        diff,
                        inner sep=1.5pt,
                    ] {};
                \end{scope}
            \end{scope}

            \node at (2,0) {$=$};

            \begin{scope}[xshift=7cm]
                \node[operator] (g) at (0,0) {$g$};
                \node[operator] (app) at (-1.5,0) {$\App$};
                \node[operator] (eta) at ($(app.140)+(-1.2,0.4)$) {$\eta$};

                \node (in1) at ($(eta.west)+(-1,0)$) {};
                \node (in2) at ($(in1 |- app.220)+(0,-0.4)$) {};
                \node (out) at ($(g.east)+(1,0)$) {};

                \draw[magenta] (in1) -- (eta);
                \draw[blue] (in2) to[out=0,in=220,in looseness=0.8] (app.220);
                \draw[blue] (eta) to[out=0,in=140] (app.140);
                \draw[blue] (app) -- (g);
                \draw[green] (g) -- (out);
                
                \begin{scope}[on background layer]
                    \node[
                        fit=(eta)(app),
                        result,
                        rotate=-20,
                        inner sep=-3pt,
                    ] {};
                    \node[
                        fit=(app)(g),
                        diff,
                        minimum height=1cm,
                        inner sep=-3pt,
                    ] {};
                \end{scope}
            \end{scope}

            \node at (2,-2.5) {$=$};

            \begin{scope}[xshift=7cm,yshift=-2.5cm]
                \node[operator] (m) at (0,0) {$m$};
                \node[operator] (g1) at ($(app |- m.140)+(0,0.4)$) {$g$};
                \node[operator] (g2) at ($(g1 |- m.220)+(0,-0.4)$) {$g$};
                \node[operator] (eta) at (eta |- g1) {$\eta$};

                \node (in1) at (in1 |- eta) {};
                \node (in2) at (in1 |- g2) {};
                \node (out) at (out |- m) {};

                \draw[magenta] (in1) -- (eta);
                \draw[blue] (in2) -- (g2);
                \draw[blue] (eta) -- (g1);
                \draw[green] (g1) to[out=0,in=140] (m.140);
                \draw[green] (g2) to[out=0,in=220] (m.220);
                \draw[green] (m) -- (out);
                
                \begin{scope}[on background layer]
                    \coordinate (F1) at ($(g1)+(-0.5,0.8)$);
                    \coordinate (F2) at ($(g2)+(-0.5,-0.8)$);
                    \coordinate (F3) at ($(m)+(1.2,0)$);
                    \draw[result,rounded corners=20pt] (F1) -- (F2) [rounded corners=25pt] -- (F3) [rounded corners=20pt] -- cycle;
                    \node[
                        fit=(eta)(g1),
                        diff,
                        minimum height=1cm,
                        inner sep=-3pt,
                    ] {};
                \end{scope}
            \end{scope}

            \node at (2,-5) {$=$};

            \begin{scope}[xshift=7cm,yshift=-5cm]
                \node[operator] (m) at (0,0) {$m$};
                \node[operator] (f) at ($(app |- m.140)+(0,0.4)$) {$f$};
                \node[operator] (g) at ($(eta |- m.220)+(0,-0.4)$) {$g$};

                \node (in1) at (in1 |- f) {};
                \node (in2) at (in1 |- g) {};
                \node (out) at (out |- m) {};

                \draw[magenta] (in1) -- (f);
                \draw[blue] (in2) -- (g);
                \draw[green] (f) to[out=0,in=140] (m.140);
                \draw[green] (g) to[out=0,in=220,in looseness=0.8] (m.220);
                \draw[green] (m) -- (out);
                
                \begin{scope}[on background layer]
                    \node[
                        fit=(f),
                        result,
                        inner sep=1.5pt,
                    ] {};
                \end{scope}
            \end{scope}
        \end{tikzpicture}
    \end{equation*}
\end{proof}












\begin{toappendix}
\begin{lemma}\label{lem:eta-app}
    \begin{equation*}
        \begin{tikzpicture}
            \begin{scope}
                \node[operator] (app) at (0,0) {$\App$};
                \node[operator] (eta) at ($(app.140)+(-1.2,0.4)$) {$\eta$};
                \node (in1) at ($(eta.west)+(-1,0)$) {};
                \node (in2) at ($(in1 |- app.220) + (0,-0.4)$) {};
                \node (out) at ($(app.east)+(1,0)$) {};

                \draw[magenta] (in1) -- (eta);
                \draw[blue] (eta) to[out=0,in=140] (app.140);
                \draw[blue] (in2) to[out=0,in=220,in looseness=0.6] (app.220);
                \draw[blue] (app) -- (out);
            \end{scope}
            \node at (2,0) {$=$};
            \begin{scope}[xshift=4.2cm]
                \node[operator,fill=consbg] (op) at (0,0) {$\cons$};
                \node (in1) at ($(op.140)+(-1.2,0.4)$) {};
                \node (in2) at ($(op.220)+(-1.2,-0.4)$) {};
                \node (out) at ($(op.east)+(1,0)$) {};
                \draw[magenta] (in1) to[out=0,in=140] (op.140);
                \draw[blue] (in2.east) to[out=0,in=220] (op.220);
                \draw[blue] (op.east) -- (out);
            \end{scope}
        \end{tikzpicture}
    \end{equation*}
\end{lemma}
\begin{proof}
    \begin{equation*}
        \begin{tikzpicture}
            \begin{scope}
                \node[operator] (app) at (0,0) {$\App$};
                \node[operator] (eta) at ($(app.140)+(-1.2,0.4)$) {$\eta$};
                \node (in1) at ($(eta.west)+(-1,0)$) {};
                \node (in2) at ($(in1 |- app.220) + (0,-0.4)$) {};
                \node (out) at ($(app.east)+(1,0)$) {};
                
                \draw[magenta] (in1) -- (eta);
                \draw[blue] (eta) to[out=0,in=140] (app.140);
                \draw[blue] (in2) to[out=0,in=220,in looseness=0.6] (app.220);
                \draw[blue] (app) -- (out);

                \begin{scope}[on background layer]
                    \node[
                        fit=(eta),
                        diff,
                        inner sep=3pt
                    ] {};
                \end{scope}
            \end{scope}
            
            \node at (2,0) {$=$};

            \begin{scope}[xshift=6cm]
                \node[operator,fill=Appbg] (app) at (0,0) {$\App$};
                \node[operator,fill=consbg] (cons) at ($(app.140)+(-1.2,0.4)$) {$\cons$};
                \node[operator] (nil) at ($(cons.220)+(-1,-0.2)$) {$\nil$};

                \node (in1) at ($(cons.140)+(-1.2,0.2)$) {};
                \node (in2) at ($(in1 |- app.220)+(0,-0.6)$) {};
                \node (out) at ($(app.east)+(1,0)$) {};

                \draw[magenta] (in1) to[out=0,in=140] (cons.140);
                \draw[blue] (nil) to[out=0,in=220] (cons.220);
                \draw[blue] (in2) to[out=0,in=220,in looseness=0.8] (app.220);
                \draw[blue] (cons) to[out=0,in=140] (app.140);
                \draw[blue] (app) -- (out);

                \begin{scope}[on background layer]
                    \node[
                        fit=(nil)(cons),
                        result,
                        rotate=20,
                        inner sep=-1.5pt,
                    ] {};
                    \node[
                        fit=(cons)(app),
                        diff,
                        rotate=-20,
                        inner sep=-1.5pt,
                    ] {};
                \end{scope}
            \end{scope}
            
            \node at (2,-2.5) {$=$};

            \begin{scope}[xshift=6cm,yshift=-2.5cm]
                \node[operator,fill=consbg] (cons) at (0,0) {$\cons$};
                \node[operator,fill=Appbg] (app) at ($(cons.220)+(-1.2,-0.4)$) {$\App$};
                \node[operator] (nil) at ($(app.140)+(-1,0.2)$) {$\nil$};

                \node (in1) at ($(in1 |- cons.140)+(0,0.6)$) {};
                \node (in2) at ($(in2 |- app.220)+(0,-0.2)$) {};
                \node (out) at (out |- cons) {};

                \draw[magenta] (in1) to[out=0,in=140,in looseness=0.8] (cons.140);
                \draw[blue] (nil) to[out=0,in=140] (app.140);
                \draw[blue] (in2) to[out=0,in=220] (app.220);
                \draw[blue] (app) to[out=0,in=220] (cons.220);
                \draw[blue] (cons) -- (out);

                \begin{scope}[on background layer]
                    \node[
                        fit=(app)(cons),
                        result,
                        rotate=20,
                        inner sep=-1.5pt,
                    ] {};
                    \node[
                        fit=(nil)(app),
                        diff,
                        rotate=-20,
                        inner sep=-1.5pt,
                    ] {};
                \end{scope}
            \end{scope}

            \node at (2,-5) {$=$};

            \begin{scope}[xshift=6cm,yshift=-5cm]
                \node[operator,fill=consbg] (cons) at (0,0) {$\cons$};
                \node[operator] (id) at ($(cons.220)+(-1.2,-0.4)$) {$\id$};

                \node (in1) at ($(in1 |- cons.140)+(0,0.6)$) {};
                \node (in2) at ($(in2 |- id.220)+(0,-0.2)$) {};
                \node (out) at (out |- cons) {};

                \draw[magenta] (in1) to[out=0,in=140,in looseness=0.8] (cons.140);
                \draw[blue] (in2) to[out=0,in=220] (id.220);
                \draw[blue] (id) to[out=0,in=220] (cons.220);
                \draw[blue] (cons) -- (out);

                \begin{scope}[on background layer]
                    \node[
                        fit=(id),
                        result,
                        inner sep=1.5pt,
                    ] {};
                \end{scope}
            \end{scope}

            \node at (2,-7.5) {$=$};

            \begin{scope}[xshift=6cm,yshift=-7.5cm]
                \node[operator,fill=consbg] (cons) at (0,0) {$\cons$};

                \node (in1) at ($(in1 |- cons.140)+(0,0.6)$) {};
                \node (in2) at ($(in2 |- cons.220)+(0,-0.6)$) {};
                \node (out) at (out |- cons) {};

                \draw[magenta] (in1) to[out=0,in=140,in looseness=0.8] (cons.140);
                \draw[blue] (in2) to[out=0,in=220,in looseness=0.8] (cons.220);
                \draw[blue] (cons) -- (out);
            \end{scope}
        \end{tikzpicture}
    \end{equation*}
\end{proof}
\end{toappendix}



\begin{theorem}[Lemma~3.4.2 of \cite{fioreListObjectsAlgebraic2017}]
  If $\cat{C}$ admits parametrised list objects for every $A$,
  then $\mathsf{Mon}(\cat{C})$ is monadic over $\cat{C}$.
\end{theorem}



\begin{theorem}[Lemma~3.4.3 of \cite{fioreListObjectsAlgebraic2017}]
  If $(\cat{C},\otimes,1)$ is cartesian, then the induced monad $T:\cat{C}\to\cat{C}$ is strong.
\end{theorem}


\section{Symmetric List Objects}

Hereafter, we assume $\cat{C}$ is symmetric strict monoidal.

\begin{definition}
  A \textit{parametrised symmetric list object} or \textit{ilst} on $A$ is a list object $L(A)$ satisfying:

  \begin{figure}[H]
    \centering

    \begin{tikzpicture}
      \begin{scope}
        \node[operator,fill=consbg] (cons2) at (0,0) {$\cons$};
        \node[anchor=east,operator,fill=consbg] (cons1) at ($(cons2.220)+(-1,-0.4)$) {$\cons$};
        \node[anchor=west] (out) at ($(cons2.east)+(1,0)$) {};
        \node[anchor=east] (in1) at ($(cons2.140)+(-1,0.6)$) {};
        \node[anchor=east] (in2) at ($(cons1.140)+(-1,0.2)$) {};
        \node[anchor=east] (in3) at ($(cons1.220)+(-1,-0.2)$) {};

        \draw[blue] (cons1.east) to[out=0,in=220] node[inputwire,pos=0.5,yshift=-2pt,below,sloped=false] {$L(A)$} (cons2.220);
        \draw[magenta] (in2.east |- in1) to[out=0,in=140] node[inputwire,pos=0.1,above,sloped=false] {$A$} (cons2.140);

        \draw[blue] (cons2.east) -- (out.west) node[outputwire,above] {$L(A)$};
        \draw[magenta] (in2.east) to[out=0,in=140] node[inputwire,above,sloped=false] {$A$} (cons1.140);
        \draw[blue] (in3.east) to[out=0,in=220] node[inputwire,below,sloped=false] {$L(A)$} (cons1.220);
      \end{scope}

      \node at (2,0) {$=$};

      \begin{scope}[xshift=5.6cm]
        \node[operator,fill=consbg] (cons2) at (0,0) {$\cons$};
        \node[anchor=east,operator,fill=consbg] (cons1) at ($(cons2.220)+(-1,-0.4)$) {$\cons$};
        \node[anchor=west] (out) at ($(cons2.east)+(1,0)$) {};
        \node[anchor=east] (in1) at ($(cons2.140)+(-1,0.6)$) {};
        \node[anchor=east] (in2) at ($(cons1.140)+(-1,0.2)$) {};
        \node[anchor=east] (in3) at ($(cons1.220)+(-1,-0.2)$) {};

        \draw[blue] (cons1.east) to[out=0,in=220] node[inputwire,pos=0.5,yshift=-2pt,below,sloped=false] {$L(A)$} (cons2.220);
        \draw[magenta] (in2.east |- in1) to[out=0,in=140] node[inputwire,pos=0.1,above,sloped=false] {$A$} (cons1.140);

        \draw[blue] (cons2.east) -- (out.west) node[outputwire,above] {$L(A)$};
        \draw[magenta] (in2.east) to[out=0,in=140] node[inputwire,pos=0.1,below,sloped=false] {$A$} (cons2.140);
        \draw[blue] (in3.east) to[out=0,in=220] node[inputwire,below,sloped=false] {$L(A)$} (cons1.220);
      \end{scope}
    \end{tikzpicture}%
    \caption*{ilst axiom}
  \end{figure}
\end{definition}



\begin{toappendix}

\begin{lemma}\label{lem:eta-app braided}
    \begin{equation*}
        \begin{tikzpicture}
            \begin{scope}
                \node[operator] (app) at (0,0) {$\App$};
                \node[operator] (eta) at ($(app.220)+(-1.2,-0.4)$) {$\eta$};
                \node (in2) at ($(eta)+(-1.5,0)$) {};
                \node (in1) at ($(in2 |- app.140) + (0,0.4)$) {};
                \node (out) at ($(app.east)+(1,0)$) {};

                \draw[magenta] (in1) to[out=0,in=180] (eta);
                \draw[blue] (in2) to[out=0,in=140,out looseness=0.5,in looseness=2] (app.140);
                \draw[blue] (eta) to[out=0,in=220] (app.220);
                \draw[blue] (app) -- (out);
            \end{scope}

            \node at (2,0) {$=$};

            \begin{scope}[xshift=4.2cm]
                \node[operator,fill=consbg] (cons) at (0,0) {$\cons$};
                \node (in1) at ($(cons.140)+(-1.2,0.4)$) {};
                \node (in2) at ($(cons.220)+(-1.2,-0.4)$) {};
                \node (out) at ($(cons.east)+(1,0)$) {};
                \draw[magenta] (in1) to[out=0,in=140] (cons.140);
                \draw[blue] (in2.east) to[out=0,in=220] (cons.220);
                \draw[blue] (cons.east) -- (out);
            \end{scope}
        \end{tikzpicture}
    \end{equation*}
\end{lemma}
\begin{proof}
    Let the two maps be $f'$ and $g'$, respectively. We claim that both maps are equal to $\iter(\eta,\cons)\circ\beta$, or equivalently, that $f:=f'\circ\beta$ and $g:=g'\circ\beta$ are the iterator $\iter(\eta,\cons)\circ\beta^2=\iter(\eta,\cons)$.
    
    $\nil$-$\beta$:
    \begin{equation*}
        \begin{tikzpicture}
            \begin{scope}[yshift=2cm]
                \node[operator] (f) at (0,0) {$f$};
                \node[anchor=east,operator] (nil) at ($(f.140)+(-1,0.4)$) {$\nil$};
                \node (in) at ($(f.220)+(-1.5,-0.4)$) {};
                \node (out) at ($(f.east)+(1,0)$) {};
                \draw[blue] (nil.east) to[out=0,in=140] (f.140);
                \draw[magenta] (in) to[out=0,in=220] (f.220);
                \draw[blue] (f.east) -- (out);
            \end{scope}

            \node at (1.55,2) {$=$};

            \begin{scope}[xshift=4.5cm,yshift=2cm]
                \node[operator] (app) at (0,0) {$\App$};
                \node[operator] (eta) at ($(app.220)+(-1.2,-0.4)$) {$\eta$};
                \node[operator] (nil) at ($(app.140) + (-2.2,0.4)$) {$\nil$};
                
                \node (in) at ($(eta.west)+(-0.8,0)$) {};
                \node (out) at ($(app.east)+(0.8,0)$) {};

                \draw[blue] (nil) to[out=0,in=140,in looseness=0.8] (app.140);
                \draw[magenta] (in) -- (eta);
                \draw[blue] (eta) to[out=0,in=220] (app.220);
                \draw[blue] (app) -- (out);

                \begin{scope}[on background layer]
                    \node[
                        fit=(nil)(app),
                        diff,
                        rotate=-12.5,
                        inner sep=-4.5pt,
                    ] {};
                \end{scope}
            \end{scope}

            \node at (1.55,0) {$=$};

            \begin{scope}[xshift=4.5cm,yshift=0.25cm]
                \node[operator] (id) at (0,0) {$\id$};
                \node[operator] (eta) at ($(id.220)+(-1.2,-0.4)$) {$\eta$};

                \node (in) at ($(eta.west)+(-0.8,0)$) {};
                \node (out) at ($(id.east)+(0.8,0)$) {};

                \draw[magenta] (in) -- (eta);
                \draw[blue] (eta) to[out=0,in=220] (id.220);
                \draw[blue] (id) -- (out);

                \begin{scope}[on background layer]
                    \node[
                        fit=(id),
                        result,
                        inner sep=1.5pt,
                    ] {};
                \end{scope}
            \end{scope}

            \node at (1.55,-2) {$=$};

            \begin{scope}[xshift=4.5cm,yshift=-2cm]
                \node[operator] (eta) at (eta |- 0,0) {$\eta$};

                \node (in) at (in |- eta) {};
                \node (out) at (out |- eta) {};

                \draw[magenta] (in) -- (eta);
                \draw[blue] (eta) -- (out);
            \end{scope}
        \end{tikzpicture}
        \qquad
        \begin{tikzpicture}
            \begin{scope}[yshift=2cm]
                \node[operator] (g) at (0,0) {$g$};
                \node[anchor=east,operator] (nil) at ($(g.140)+(-1,0.4)$) {$\nil$};
                \node (in) at ($(g.220)+(-1.5,-0.4)$) {};
                \node (out) at ($(g.east)+(1,0)$) {};
                \draw[blue] (nil.east) to[out=0,in=140] (g.140);
                \draw[magenta] (in) to[out=0,in=220] (g.220);
                \draw[blue] (g.east) -- (out);
            \end{scope}

            \node at (1.6,2) {$=$};

            \begin{scope}[xshift=4.3cm,yshift=2cm]
                \node[operator,fill=consbg] (cons) at (0,0) {$\cons$};
                \node[operator] (nil) at ($(cons.140)+(-2,0.4)$) {$\nil$};

                \node (in) at ($(cons.220)+(-1.9,-0.4)$) {};
                \node (out) at ($(cons.east)+(0.8,0)$) {};
                
                \draw[blue] (nil) to[out=0,in=220,in looseness=2.5] (cons.220);
                \draw[magenta] (in) to[out=0,in=140,out looseness=1.3,in looseness=2.5] (cons.140);
                \draw[blue] (cons) -- (out);
            \end{scope}

            \node at (1.6,0) {$=$};

            \begin{scope}[xshift=4.3cm]
                \node[operator,fill=consbg] (cons) at (0,0) {$\cons$};
                \node[operator] (nil) at ($(cons.220)+(-0.8,-0.4)$) {$\nil$};

                \node (in) at ($(cons.220)+(-1.9,-0.4)$) {};
                \node (out) at ($(cons.east)+(0.8,0)$) {};
                
                \draw[blue] (nil) to[out=0,in=220] (cons.220);
                \draw[magenta] (in) to[out=0,in=140,out looseness=1.3,in looseness=2.5] (cons.140);
                \draw[blue] (cons) -- (out);

                \begin{scope}[on background layer]
                    \node[
                        fit=(nil)(cons),
                        diff,
                        rotate=30,
                        inner sep=-3pt,
                    ] {};
                \end{scope}
            \end{scope}

            \node at (1.6,-2) {$=$};

            \begin{scope}[xshift=4.3cm,yshift=-2cm]
                \node[operator] (eta) at ($(nil |- 0,0)+(0.6,0)$) {$\eta$};

                \node (in) at (in |- eta) {};
                \node (out) at (out |- eta) {};

                \draw[magenta] (in) -- (eta);
                \draw[blue] (eta) -- (out);
            \end{scope}
        \end{tikzpicture}
    \end{equation*}



    $\cons$-$\beta$:

    \begin{equation*}
        \begin{tikzpicture}[baseline=(base),scale=0.8]
            \node (base) at (0,0) {};
            \begin{scope}
                \node[operator] (f) at (0,0) {$f$};
                \node[operator,fill=consbg] (cons) at ($(f.140)+(-1.5,0.4)$) {$\cons$};

                \node (in1) at ($(cons.140)+(-1,0.2)$) {};
                \node (in2) at ($(cons.220)+(-1,-0.2)$) {};
                \node (in3) at ($(in1 |- f.220)+(0,-0.6)$) {};
                \node (out) at ($(f.east)+(1,0)$) {};

                \draw[magenta] (in1) to[out=0,in=140] (cons.140);
                \draw[blue] (in2) to[out=0,in=220] (cons.220);
                \draw[magenta] (in3) to[out=0,in=220] (f.220);
                \draw[blue] (cons) to[out=0,in=140] (f.140);
                \draw[blue] (f) -- (out);
            \end{scope}

            \node at (1.7,0) {$=$};

            \begin{scope}[xshift=4.9cm]
                \node[operator] (app) at (0,0) {$\App$};
                \node[operator,fill=consbg] (cons) at ($(app.140)+(-1.2,0.4)$) {$\cons$};
                \node[operator] (eta) at ($(app.220)+(-1.2,-0.4)$) {$\eta$};

                \node (in1) at ($(cons.140)+(-1,0.2)$) {};
                \node (in2) at ($(cons.220)+(-1,-0.2)$) {};
                \node (in3) at (in1 |- eta) {};
                \node (out) at ($(app.east)+(1,0)$) {};

                \draw[magenta] (in1) to[out=0,in=140] (cons.140);
                \draw[blue] (in2) to[out=0,in=220] (cons.220);
                \draw[magenta] (in3) -- (eta);
                \draw[blue] (cons) to[out=0,in=140] (app.140);
                \draw[blue] (eta) to[out=0,in=220] (app.220);
                \draw[blue] (app) -- (out);

                \begin{scope}[on background layer]
                    \node[
                        fit=(cons)(app),
                        diff,
                        rotate=-20,
                        inner sep=-3pt,
                    ] {};
                \end{scope}
            \end{scope}

            \node at (1.7,-2.5) {$=$};

            \begin{scope}[xshift=4.9cm,yshift=-3cm]
                \node[operator] (app) at (0,0) {$\App$};
                \node[operator,fill=consbg] (cons) at ($(app.40)+(1.2,0.4)$) {$\cons$};
                \node[operator] (eta) at ($(app.220)+(-1.2,-0.4)$) {$\eta$};

                \node (in1) at ($(in1 |- cons.140)+(0,0.6)$) {};
                \node (in2) at ($(in1 |- app.140)+(0,0.4)$) {};
                \node (in3) at (in1 |- eta) {};
                \node (out) at ($(cons.east)+(1,0)$) {};

                \draw[magenta] (in1) to[out=0,in=140,in looseness=0.8] (cons.140);
                \draw[blue] (in2) to[out=0,in=140,in looseness=0.8] (app.140);
                \draw[magenta] (in3) -- (eta);
                \draw[blue] (app) to[out=0,in=220] (cons.220);
                \draw[blue] (eta) to[out=0,in=220] (app.220);
                \draw[blue] (cons) -- (out);
            \end{scope}

            \node at (1.7,-5.5) {$=$};

            \begin{scope}[xshift=4.9cm,yshift=-5.5cm]
                \node[operator,fill=consbg] (cons) at (cons |- 0,0) {$\cons$};
                \node[operator] (f) at ($(cons.220)+(-1.7,-0.4)$) {$f$};

                \node (in1) at ($(in1 |- cons.140)+(0,0.6)$) {};
                \node (in2) at ($(in1 |- f.140)+(0,0.4)$) {};
                \node (in3) at ($(in1 |- f.220)+(0,-0.4)$) {};
                \node (out) at ($(cons.east)+(1,0)$) {};

                \draw[magenta] (in1) to[out=0,in=140,in looseness=0.8] (cons.140);
                \draw[blue] (in2) to[out=0,in=140,in looseness=0.8] (f.140);
                \draw[magenta] (in3) to[out=0,in=220] (f.220);
                \draw[blue] (f) to[out=0,in=220] (cons.220);
                \draw[blue] (cons) -- (out);
            \end{scope}
        \end{tikzpicture}
        \qquad
        \begin{tikzpicture}[baseline=(base),scale=0.8]
            \node (base) at (0,0) {};
            \begin{scope}
                \node[operator] (g) at (0,0) {$g$};
                \node[operator,fill=consbg] (cons) at ($(g.140)+(-1.5,0.4)$) {$\cons$};

                \node (in1) at ($(cons.140)+(-1,0.2)$) {};
                \node (in2) at ($(cons.220)+(-1,-0.2)$) {};
                \node (in3) at ($(in1 |- g.220)+(0,-0.6)$) {};
                \node (out) at ($(g.east)+(1,0)$) {};

                \draw[magenta] (in1) to[out=0,in=140] (cons.140);
                \draw[blue] (in2) to[out=0,in=220] (cons.220);
                \draw[magenta] (in3) to[out=0,in=220] (g.220);
                \draw[blue] (cons) to[out=0,in=140] (g.140);
                \draw[blue] (g) -- (out);
            \end{scope}

            \node at (1.7,0) {$=$};

            \begin{scope}[xshift=5.1cm]
                \node[operator,fill=consbg] (cons2) at (0,0) {$\cons$};
                \node[operator,fill=consbg] (cons1) at ($(cons2.140)+(-1.5,0.4)$) {$\cons$};

                \node (in1) at ($(cons1.140)+(-1,0.2)$) {};
                \node (in2) at ($(cons1.220)+(-1,-0.2)$) {};
                \node (in3) at ($(in1 |- cons2.220)+(0,-0.6)$) {};
                \node (out) at ($(cons2.east)+(1,0)$) {};

                \draw[magenta] (in1) to[out=0,in=140] (cons1.140);
                \draw[blue] (in2) to[out=0,in=220] (cons1.220);
                \draw[magenta] (in3) to[out=0,in=140,out looseness=2,in looseness=1.2] (cons2.140);
                \draw[blue] (cons1) to[out=0,in=220,in looseness=2.5] (cons2.220);
                \draw[blue] (cons2) -- (out);   
            \end{scope}

            \node at (1.7,-2.5) {$=$};

            \begin{scope}[xshift=5.1cm,yshift=-2.5cm]
                \node[operator,fill=consbg] (cons2) at (0,0) {$\cons$};
                \node[operator,fill=consbg] (cons1) at ($(cons2.220)+(-0.8,-0.4)$) {$\cons$};
                \node (anchor) at ($(cons2.140)+(-1.2,0.4)$) {};

                \node (in1) at ($(in1 |- anchor)+(0,0.4)$) {};
                \node (in2) at ($(in1 |- anchor)+(0,-0.4)$) {};
                \node (in3) at ($(in1 |- cons2.220)+(0,-0.6)$) {};
                \node (out) at ($(cons2.east)+(1,0)$) {};

                \draw[magenta] (in1) to[out=0,in=140] (cons1.140);
                \draw[blue] (in2) to[out=0,in=220] (cons1.220);
                \draw[magenta] (in3) to[out=0,in=140,out looseness=0.8,in looseness=3] (cons2.140);
                \draw[blue] (cons1) to[out=0,in=220] (cons2.220);
                \draw[blue] (cons2) -- (out);

                \begin{scope}[on background layer]
                    \node[
                        fit=(cons1)(cons2),
                        diff,
                        rotate=25,
                        inner sep=-3pt,
                    ] {};
                \end{scope}
            \end{scope}

            \node at (1.7,-5) {$=$};

            \begin{scope}[xshift=5.1cm,yshift=-5cm]
                \node[operator,fill=consbg] (cons2) at (0,0) {$\cons$};
                \node[operator,fill=consbg] (cons1) at ($(cons2.220)+(-0.8,-0.4)$) {$\cons$};
                \node (anchor) at ($(cons2.140)+(-1.2,0.4)$) {};

                \node (in1) at ($(in1 |- anchor)+(0,0.2)$) {};
                \node (in2) at ($(in1 |- cons1.140)+(0,0.4)$) {};
                \node (in3) at ($(in1 |- cons1.220)+(0,-0.4)$) {};
                \node (out) at ($(cons2.east)+(1,0)$) {};

                \draw[magenta] (in1) to[out=0,in=140] (cons2.140);
                \draw[blue] (in2) to[out=0,in=220,in looseness=2.5] (cons1.220);
                \draw[magenta] (in3) to[out=0,in=140,in looseness=2.5] (cons1.140);
                \draw[blue] (cons1) to[out=0,in=220] (cons2.220);
                \draw[blue] (cons2) -- (out);

                \begin{scope}[on background layer]
                    \node[
                        fit=(cons1)(cons2),
                        result,
                        rotate=25,
                        inner sep=-3pt,
                    ] {};
                \end{scope}
            \end{scope}

            \node at (1.7,-8) {$=$};

            \begin{scope}[xshift=5.1cm,yshift=-8cm]
                \node[operator,fill=consbg] (cons) at (0,0) {$\cons$};
                \node[operator] (g) at ($(cons.220)+(-1,-0.4)$) {$g$};
                \node (anchor) at ($(cons.140)+(-1.2,0.4)$) {};

                \node (in1) at ($(in1 |- anchor)+(0,0.2)$) {};
                \node (in2) at ($(in1 |- g.140)+(0,0.4)$) {};
                \node (in3) at ($(in1 |- g.220)+(0,-0.4)$) {};
                \node (out) at ($(cons.east)+(1,0)$) {};

                \draw[magenta] (in1) to[out=0,in=140] (cons.140);
                \draw[blue] (in2) to[out=0,in=140] (g.140);
                \draw[magenta] (in3) to[out=0,in=220] (g.220);
                \draw[blue] (g) to[out=0,in=220] (cons.220);
                \draw[blue] (cons) -- (out);
            \end{scope}
        \end{tikzpicture}
    \end{equation*}
\end{proof}

\end{toappendix}




\begin{theoremrep}
  Let $(L(A),\nil,\cons)$ be a list object on an object $A$. Then, the following are equivalent:
  \begin{enumerate}
    \item $(L(A),\nil,\cons,\iter)$ is a symmetric list object;
    \item $(L(A),\nil,\App)$ is a commutative monoid object;
    \item $(L(A),\nil,\App)$ is the free commutative monoid object on $A$;
  \end{enumerate}
\end{theoremrep}

\begin{proof}
    $(1\Rightarrow2)$ We have already shown that list objects are monoids, so we need only show that $\App$ is commutative:
    \begin{equation*}
        \begin{tikzpicture}
            \begin{scope}
                \node[operator,fill=Appbg] (app) at (0,0) {$\App$};
                
                \node (in1) at ($(app.140)+(-1.5,0.4)$) {};
                \node (in2) at ($(app.220)+(-1.5,-0.4)$) {};
                \node (out) at ($(app.east)+(0.8,0)$) {};
            
                \draw[blue] (in1) to[out=0,in=140] (app.140);
                \draw[blue] (in2) to[out=0,in=220] (app.220);
                \draw[blue] (app) -- (out);
            \end{scope}

            \node at (2,0) {$=$};

            \begin{scope}[xshift=5cm]
                \node[operator,fill=Appbg] (app) at (0,0) {$\App$};
                
                \node (in1) at ($(app.140)+(-2,0.4)$) {};
                \node (in2) at ($(app.220)+(-2,-0.4)$) {};
                \node (out) at ($(app.east)+(0.8,0)$) {};
            
                \draw[blue] (in1) to[out=0,in=220,out looseness=1.3,in looseness=2.5] (app.220);
                \draw[blue] (in2) to[out=0,in=140,out looseness=1.3,in looseness=2.5] (app.140);
                \draw[blue] (app) -- (out);
            \end{scope}
        \end{tikzpicture}
    \end{equation*}





    $(2\Rightarrow1)$ $L(A)$ is a list object by assumption, so it suffices to show the ilst axiom only.

    \begin{equation*}
        \begin{tikzpicture}
            \begin{scope}
                \node[operator,fill=consbg] (cons2) at (0,0) {$\cons$};
                \node[anchor=east,operator,fill=consbg] (cons1) at ($(cons2.220)+(-1,-0.4)$) {$\cons$};
                
                \node (in1) at ($(cons2.140)+(-1,0.6)$) {};
                \node (in2) at ($(cons1.140)+(-1,0.2)$) {};
                \node (in3) at ($(cons1.220)+(-1,-0.2)$) {};
                \node (out) at ($(cons2.east)+(1,0)$) {};

                \draw[blue] (cons1.east) to[out=0,in=220] (cons2.220);
                \draw[magenta] (in2.east |- in1) to[out=0,in=140] (cons2.140);

                \draw[blue] (cons2.east) -- (out.west);
                \draw[magenta] (in2.east) to[out=0,in=140] (cons1.140);
                \draw[blue] (in3.east) to[out=0,in=220] (cons1.220);

                \begin{scope}[on background layer]
                    \node [
                        fit=(cons1),
                        diff,
                        inner sep=1.5pt
                    ] {};
                    \node [
                        fit=(cons2),
                        diff,
                        inner sep=1.5pt
                    ] {};
                \end{scope}
            \end{scope}
        
            \node at (2,0) {$=$};
        
            \begin{scope}[xshift=7cm]
                \node[operator,fill=consbg] (app2) at (0,0) {$\App$};
                \node[operator,fill=consbg] (app1) at ($(app2.220)+(-1.5,-0.4)$) {$\App$};
                \node[operator] (eta1) at ($(app1.140)+(-1.5,0.4)$) {$\eta$};
                \node[operator] (eta2) at ($(app2.140)+(-1.5,0.4)$) {$\eta$};
                
                \node (in2) at ($(eta1)+(-1,0)$) {};
                \node (in1) at (in2 |- eta2) {};
                \node (in3) at ($(app1.220)+(-2.5,-0.4)$) {};
                \node (out) at ($(app2.east)+(1,0)$) {};

                \draw[magenta] (in1) -- (eta2);
                \draw[magenta] (in2) -- (eta1);
                \draw[blue] (in3) to[out=0,in=220] (app1.220);
                \draw[blue] (eta1) to[out=0,in=140] (app1.140);
                \draw[blue] (eta2) to[out=0,in=140] (app2.140);
                \draw[blue] (app1) to[out=0,in=220] (app2.220);
                \draw[blue] (app2) -- (out);

                \begin{scope}[on background layer]
                    \node [
                        fit=(eta1)(app1),
                        result,
                        rotate=-20,
                        inner sep=-3pt
                    ] {};
                    \node [
                        fit=(eta2)(app2),
                        result,
                        rotate=-20,
                        inner sep=-3pt
                    ] {};
                    \node [
                        fit=(app1)(app2),
                        diff,
                        rotate=15,
                        inner sep=-3pt
                    ] {};
                \end{scope}
            \end{scope}
                
            \node at (2,-3) {$=$};
        
            \begin{scope}[xshift=7cm,yshift=-3cm]
                \node[operator,fill=consbg] (app2) at (0,0) {$\App$};
                \node[operator,fill=consbg] (app1) at ($(app2.140)+(-1.5,0.4)$) {$\App$};
                \node[operator] (eta1) at ($(app1.140)+(-1.5,0.4)$) {$\eta$};
                \node[operator] (eta2) at ($(app1.220)+(-1.5,-0.4)$) {$\eta$};
                
                \node (in1) at (in1 |- eta1) {};
                \node (in2) at (in2 |- eta2) {};
                \node (in3) at ($(in3 |- app2.220)+(0,-0.6)$) {};
                \node (out) at (out |- app2) {};

                \draw[magenta] (in1) -- (eta1);
                \draw[magenta] (in2) -- (eta2);
                \draw[blue] (in3) to[out=0,in=220,in looseness=0.8] (app2.220);
                \draw[blue] (eta1) to[out=0,in=140] (app1.140);
                \draw[blue] (eta2) to[out=0,in=220] (app1.220);
                \draw[blue] (app1) to[out=0,in=140] (app2.140);
                \draw[blue] (app2) -- (out);

                \begin{scope}[on background layer]
                    \coordinate (F1) at ($(eta1)+(-0.5,0.7)$);
                    \coordinate (F2) at ($(eta2)+(-0.5,-0.7)$);
                    \coordinate (F3) at ($(app1)+(1.2,0)$);
                    \draw[diff,rounded corners=15pt] (F1) -- (F2) [rounded corners=30pt] -- (F3) [rounded corners=15pt] -- cycle;
                    \node [
                        fit=(app1)(app2),
                        result,
                        rotate=-15,
                        inner sep=-3pt
                    ] {};
                \end{scope}
            \end{scope}

            \node at (2,-6) {$=$};
        
            \begin{scope}[xshift=7cm,yshift=-6cm]
                \node[operator,fill=consbg] (app2) at (0,0) {$\App$};
                \node[operator,fill=consbg] (app1) at ($(app2.140)+(-1.5,0.4)$) {$\App$};
                \node[operator] (eta1) at ($(app1.140)+(-1.5,0.4)$) {$\eta$};
                \node[operator] (eta2) at ($(app1.220)+(-1.5,-0.4)$) {$\eta$};
                
                \node (in1) at (in1 |- eta1) {};
                \node (in2) at (in2 |- eta2) {};
                \node (in3) at ($(in3 |- app2.220)+(0,-0.6)$) {};
                \node (out) at (out |- app2) {};

                \draw[magenta] (in1) -- (eta1);
                \draw[magenta] (in2) -- (eta2);
                \draw[blue] (in3) to[out=0,in=220,in looseness=0.8] (app2.220);
                \draw[blue] (eta1) to[out=0,in=220,in looseness=2] (app1.220);
                \draw[blue] (eta2) to[out=0,in=140,in looseness=2] (app1.140);
                \draw[blue] (app1) to[out=0,in=140] (app2.140);
                \draw[blue] (app2) -- (out);

                \begin{scope}[on background layer]
                    \coordinate (F1) at ($(eta1)+(-0.5,0.7)$);
                    \coordinate (F2) at ($(eta2)+(-0.5,-0.7)$);
                    \coordinate (F3) at ($(app1)+(1.2,0)$);
                    \draw[result,rounded corners=15pt] (F1) -- (F2) [rounded corners=30pt] -- (F3) [rounded corners=15pt] -- cycle;
                \end{scope}
            \end{scope}

            \node at (2,-9) {$=$};

            \begin{scope}[xshift=7cm,yshift=-9cm]
                \node[operator,fill=consbg] (app2) at (0,0) {$\App$};
                \node[operator,fill=consbg] (app1) at ($(app2.140)+(-1.5,0.4)$) {$\App$};
                \node[operator] (eta1) at ($(app1.140)+(-1,0.4)$) {$\eta$};
                \node[operator] (eta2) at ($(app1.220)+(-1,-0.4)$) {$\eta$};
                
                \node (in1) at (in1 |- eta1) {};
                \node (in2) at (in2 |- eta2) {};
                \node (in3) at ($(in3 |- app2.220)+(0,-0.6)$) {};
                \node (out) at (out |- app2) {};

                \draw[magenta] (in1) to[out=0,in=180] (eta2);
                \draw[magenta] (in2) to[out=0,in=180] (eta1);
                \draw[blue] (in3) to[out=0,in=220,in looseness=0.8] (app2.220);
                \draw[blue] (eta1) to[out=0,in=140] (app1.140);
                \draw[blue] (eta2) to[out=0,in=220] (app1.220);
                \draw[blue] (app1) to[out=0,in=140] (app2.140);
                \draw[blue] (app2) -- (out);

                \begin{scope}[on background layer]
                    \node [
                        fit=(app1)(app2),
                        diff,
                        rotate=-15,
                        inner sep=-3pt
                    ] {};
                \end{scope}
            \end{scope}

            \node at (2,-12) {$=$};

            \begin{scope}[xshift=7cm,yshift=-12cm]
                \node[operator,fill=consbg] (app2) at (0,0) {$\App$};
                \node[operator,fill=consbg] (app1) at ($(app2.220)+(-1.5,-0.4)$) {$\App$};
                \node[operator] (eta1) at ($(eta1 |- app2.140)+(0,0.8)$) {$\eta$};
                \node[operator] (eta2) at ($(eta2 |- app1.140)+(0,0.4)$) {$\eta$};
                
                \node (in1) at (in1 |- eta1) {};
                \node (in2) at (in2 |- eta2) {};
                \node (in3) at ($(in3 |- app1.220)+(0,-0.4)$) {};
                \node (out) at (out |- app2) {};

                \draw[magenta] (in1) to[out=0,in=180] (eta2);
                \draw[magenta] (in2) to[out=0,in=180] (eta1);
                \draw[blue] (in3) to[out=0,in=220] (app1.220);
                \draw[blue] (eta1) to[out=0,in=140] (app2.140);
                \draw[blue] (eta2) to[out=0,in=140] (app1.140);
                \draw[blue] (app1) to[out=0,in=220] (app2.220);
                \draw[blue] (app2) -- (out);

                \begin{scope}[on background layer]
                    \node [
                        fit=(eta1)(app2),
                        diff,
                        rotate=-20,
                        inner sep=-6pt
                    ] {};
                    \node [
                        fit=(eta2)(app1),
                        diff,
                        rotate=-20,
                        inner sep=-3pt
                    ] {};
                \end{scope}
            \end{scope}

            \node at (2,-15) {$=$};

            \begin{scope}[xshift=5.6cm,yshift=-15cm]
                \node[operator,fill=consbg] (cons1) at (1,0) {$\cons$};
                \node[operator,fill=consbg] (cons2) at ($(cons1)+(-1.5,-0.4)$) {$\cons$};

                \node (in1) at ($(in1 |- cons1.140)+(0,0.6)$) {};
                \node (in2) at ($(in2 |- cons2.140)+(0,0.2)$) {};
                \node (in3) at ($(in3 |- cons2.220)+(0,-0.4)$) {};
                \node (out) at (out |- cons1) {};

                \draw[blue] (cons2.east) to[out=0,in=220] (cons1.220);
                \draw[magenta] (in2.east |- in1) to[out=0,in=140] (cons2.140);

                \draw[blue] (cons1.east) -- (out.west);
                \draw[magenta] (in2.east) to[out=0,in=140]  (cons1.140);
                \draw[blue] (in3.east) to[out=0,in=220] (cons2.220);
            \end{scope}
        \end{tikzpicture}
    \end{equation*}
\end{proof}

\begin{theorem}
  If $\cat{C}$ admits parametrised symmetric list objects for every $A$,
  then the forgetful functor $U:\mathsf{CMon}(\cat{C})\to\cat{C}$ is monadic.
\end{theorem}

\begin{theorem}
  If $(\cat{C},\otimes,1)$ is cartesian, then the induced monad $T:\cat{C}\to\cat{C}$ is strong, and further, commutative.
\end{theorem}

Symmetric lists give an alternative way of constructing free commutative monoids in a symmetric monoidal category which is relevant to the study of the exponential modality in linear logic. In the future, we hope to compare this to other approaches, such as in \cite{melliesExplicitFormulaFree2018}.


\clearpage

\printbibliography

\end{document}
